\documentclass[10pt,a4paper]{article}
\usepackage[margin=2.2cm]{geometry}
\usepackage{fontspec}
\setmonofont{DejaVu Sans Mono}[Scale=0.84]
\newfontfamily\unifont{DejaVu Serif}[Scale=0.9]
\usepackage{amsmath,amssymb,mathtools}
\setcounter{MaxMatrixCols}{20}
\usepackage{longtable,booktabs,array}
\usepackage{fancyvrb}
\usepackage[dvipsnames]{xcolor}
\usepackage[colorlinks=true,linkcolor=NavyBlue,urlcolor=NavyBlue]{hyperref}
\usepackage{microtype}
\usepackage{needspace}
\usepackage{tocloft}\setlength{\cftbeforesecskip}{0.45em}
\setlength{\parindent}{0pt}\setlength{\parskip}{0.55em}
\setlength{\emergencystretch}{8em}
\renewcommand{\arraystretch}{1.15}
\allowdisplaybreaks
\begin{document}
\title{Gravity on the 7+1 spacetime and the coupling of dirac16complex to it\\[0.4em]\large Einstein and Lovelock tensors, generalized frame matrices, the canonical spin connection with a vanishing total covariant derivative, the covariant Dirac equation of dirac16complex, and the primordial, modern and far-future epochs — sections 11–14 of the notebook 7+1.nb}
\author{Repository: github.com/once-ere/7plus1 — prepared for Nash\\ Notebook: 7+1.nb (Wolfram Language 15.0.1), evaluated Tue 27 Oct 2026 14:15:35\\ Document generated by tools/make\_docs.py on 2026-10-27}
\date{}
\maketitle
\begin{abstract}
This document puts the 16-component field dirac16complex of the 7+1 project into a gravitational field. The spacetime has ONE time coordinate $x_0$ and SEVEN space coordinates $x_1,\dots,x_7$; the four auxiliary space-like coordinates $x_4,\dots,x_7$ may expand or contract at their own rate. It defines and checks the curvature of a general metric, the Lovelock tensors of 7+1 spacetime (Einstein, Gauss–Bonnet, cubic; the quartic one vanishes identically in eight dimensions) and the Einstein–Lovelock field equations; the generalized frame matrices $\Gamma_\mu(X)=e^A{}_\mu(X)\,\Gamma_A$ built from the real 16×16 Dirac matrices; the canonical spin connection, for which the total covariant derivative of every frame matrix vanishes; the covariant Dirac equation, in which that spin connection couples gravity to dirac16complex; and the three cosmological epochs — primordial (near $x_0=0$), modern, far future — with the creation of dirac16complex pairs in each. The ideas follow the repository Dirac\_claude, rebuilt for 7+1: Dirac\_claude uses a 4+4 spacetime with four time-like coordinates and a warped metric, none of which is used here. Sections 11–14 of 7+1.nb contain 97 checks (97 True); the independent Python program tools/ref\_gravity.py contains 99 (99 True).
\end{abstract}
\tableofcontents
\Needspace{12\baselineskip}\section*{1. Results at a glance}\addcontentsline{toc}{section}{1. Results at a glance}
\textbf{Spacetime.} Coordinates $X=(x_0,x_1,\dots,x_7)$ with $x_0$ the time (the direction of the octonion unit $e_0=1$), $x_1,x_2,x_3$ ordinary space and $x_4,\dots,x_7$ auxiliary space. The metric $g_{\mu\nu}(X)$ has signature $(+,-,-,-,-,-,-,-)$. The frame metric and the eight real Dirac matrices are
\begin{equation*}
\eta=\mathrm{diag}(+1,-1,\dots,-1),\qquad \Gamma_0=\begin{pmatrix}0&1_8\\1_8&0\end{pmatrix},\quad \Gamma_i=\begin{pmatrix}\tau^{O}_i&0\\0&-\tau^{O}_i\end{pmatrix},\qquad \Gamma_A\Gamma_B+\Gamma_B\Gamma_A=2\eta_{AB}1_{16},
\end{equation*}
with the real antisymmetric 8×8 octonion matrices $\tau^{O}_i$ ($i=1,\dots,7$, $(\tau^O_i)^2=-1_8$, mutually anticommuting).

\textbf{Lovelock tensors} of order $k$, with the generalized Kronecker delta $\delta^{\mu_1\cdots\mu_n}_{\nu_1\cdots\nu_n}=\det[\delta^{\mu_i}_{\nu_j}]$:
\begin{gather*}
E^{(k)\mu}{}_{\nu}=-\frac{1}{2^{k+1}}\,\delta^{\mu\alpha_1\beta_1\cdots\alpha_k\beta_k}_{\nu\gamma_1\delta_1\cdots\gamma_k\delta_k}\,R^{\gamma_1\delta_1}{}_{\alpha_1\beta_1}\cdots R^{\gamma_k\delta_k}{}_{\alpha_k\beta_k},\\ E^{(0)}=-\tfrac12\,\delta,\quad E^{(1)}=G\ \text{(Einstein)},\quad E^{(4)}\equiv0\ \text{in }D=8.
\end{gather*}
\textbf{Einstein–Lovelock field equations} ($\alpha_0=2\Lambda$, $\alpha_1=1$, $\kappa=8\pi G_8$):
\begin{equation*}
\sum_{k=0}^{3}\alpha_k\,E^{(k)\mu}{}_{\nu}=\kappa\,T^{\mu}{}_{\nu}.
\end{equation*}
\textbf{Generalized frame matrices, canonical spin connection, total covariant derivative:}
\begin{align*}
\Gamma_\mu(X)&=e^A{}_\mu(X)\,\Gamma_A,\qquad \Gamma_\mu\Gamma_\nu+\Gamma_\nu\Gamma_\mu=2g_{\mu\nu}(X)\,1_{16}, \\
\omega_\mu{}^A{}_B&=e^A{}_\nu\big(\partial_\mu e_B{}^\nu+\{{}^{\nu}_{\mu\lambda}\}\,e_B{}^\lambda\big),\qquad \Omega_\mu=\tfrac14\,\omega_{\mu AB}\,\Gamma^A\Gamma^B, \\
\nabla_\mu\Gamma_\nu&=\partial_\mu\Gamma_\nu-\{{}^{\lambda}_{\mu\nu}\}\Gamma_\lambda+[\Omega_\mu,\Gamma_\nu]=0 .
\end{align*}
\textbf{Covariant Dirac equation of dirac16complex} ($D_\mu=\partial_\mu+\Omega_\mu$, $\bar\psi=\psi^\dagger\Gamma_0$):
\begin{equation*}
\big(i\,\Gamma^\mu(X)\,D_\mu-m\big)\psi=0,\qquad (\Gamma^\mu D_\mu)^2=g^{\mu\nu}\nabla_\mu D_\nu-\tfrac{R}{4}\,1_{16}.
\end{equation*}
\textbf{Three epochs} of the cosmological metric $ds^2=dx_0^2-a^2(dx_1^2+dx_2^2+dx_3^2)-b^2(dx_4^2+\cdots+dx_7^2)$: primordial Kasner solutions near $x_0=0$ (Einstein: $a=x_0^{p}$, $b=x_0^{q}$ with $p=(1\pm2\sqrt2)/7$, $q=(2\mp3\sqrt2)/14$; Gauss–Bonnet and cubic Lovelock exponents below), a modern epoch with a static auxiliary space ($b=b_0$, ordinary Friedmann cosmology with $G_4=G_8/V_\perp$), and de Sitter far futures, in which dirac16complex pairs are created with the Fermi–Dirac distribution at the Gibbons–Hawking temperature, $|\beta|^2=1/(e^{2\pi M/H}+1)$.

\Needspace{12\baselineskip}\section*{2. Conventions}\addcontentsline{toc}{section}{2. Conventions}
\begin{itemize}
\item Units: $\hbar=c=1$. A prime is $d/dx_0$.
\item Indices: Greek $\mu,\nu,\lambda,\rho,\sigma,\alpha,\beta,\gamma,\delta\in\lbrace0,\dots,7\rbrace$ are coordinate indices; Latin capitals $A,B,C\in\lbrace0,\dots,7\rbrace$ are frame indices, raised and lowered with $\eta$; repeated indices are summed. In Mathematica all lists are 1-based: $g_{\mu\nu}$ is \texttt{g[[μ+1, ν+1]]}.
\item Christoffel symbols are written $\{{}^{\lambda}_{\mu\nu}\}$ — never $\Gamma$, which in this project always denotes a Dirac matrix:
\end{itemize}
\begin{equation*}
\{{}^{\lambda}_{\mu\nu}\}=\tfrac12\,g^{\lambda\rho}\big(\partial_\mu g_{\rho\nu}+\partial_\nu g_{\rho\mu}-\partial_\rho g_{\mu\nu}\big).
\end{equation*}
Curvature (the convention of Landau and Lifshitz):
\begin{align*}
R^{\rho}{}_{\sigma\mu\nu}&=\partial_\mu\{{}^{\rho}_{\sigma\nu}\}-\partial_\nu\{{}^{\rho}_{\sigma\mu}\}+\{{}^{\rho}_{\mu\lambda}\}\{{}^{\lambda}_{\sigma\nu}\}-\{{}^{\rho}_{\nu\lambda}\}\{{}^{\lambda}_{\sigma\mu}\}, \\
R_{\sigma\nu}&=R^{\rho}{}_{\sigma\rho\nu},\qquad R=g^{\sigma\nu}R_{\sigma\nu},\qquad G^{\mu}{}_{\nu}=R^{\mu}{}_{\nu}-\tfrac12\delta^\mu_\nu R,\qquad R^{\gamma\delta}{}_{\alpha\beta}=g^{\delta\sigma}R^{\gamma}{}_{\sigma\alpha\beta}.
\end{align*}
With these conventions and the signature $(+,-,\dots,-)$ the field equations of general relativity read $G^\mu{}_\nu-\Lambda\delta^\mu_\nu=\kappa T^\mu{}_\nu$ with $T^0{}_0>0$ the energy density, and a de Sitter space has $R<0$.

\Needspace{12\baselineskip}\section*{3. Curvature of a 7+1 metric (section 11 of 7+1.nb)}\addcontentsline{toc}{section}{3. Curvature of a 7+1 metric (section 11 of 7+1.nb)}
The notebook function \texttt{curvature[g, simp]} returns, for an 8×8 metric in the coordinates $X$, the inverse metric, the Christoffel symbols, $R^\rho{}_{\sigma\mu\nu}$, $R^{\gamma\delta}{}_{\alpha\beta}$, the Ricci tensor, $R$ and $G^\mu{}_\nu$. Three test metrics are used throughout:

\begin{enumerate}
\item the \textbf{cosmological metric} $g=\mathrm{diag}(1,-a^2,-a^2,-a^2,-b^2,-b^2,-b^2,-b^2)$ with arbitrary $a(x_0)$ (\texttt{a3[x0]}) and $b(x_0)$ (\texttt{b4[x0]});
\item an \textbf{inhomogeneous metric} $g=\mathrm{diag}(n^2,-a^2,-a^2,-a^2,-b^2,-b^2,-b^2,-b^2)$ in which $n,a,b$ are explicit polynomials of BOTH $x_0$ and $x_4$: $n=1+x_0x_4/5+x_0^2/11$, $a=1+x_0/3+x_4^2/7$, $b=2-x_0x_4/4+x_0^2/9+x_4/13$ (\texttt{nP}, \texttt{aP}, \texttt{bP}); its identities are checked exactly (rational numbers) at the point $x_0=1/2$, $x_4=1/3$;
\item a \textbf{plane gravitational wave} (pp-wave) $ds^2=dx_0^2-dx_1^2-\cdots-dx_7^2+F(x_0-x_1,x_2,x_3)\,(dx_0-dx_1)^2$ with an arbitrary profile $F$ (\texttt{fPP}), which is not diagonal.
\end{enumerate}
For the cosmological metric, with the Hubble rates $H_a=a'/a$ and $H_b=b'/b$, the Riemann tensor in the orthonormal frame is diagonal on index pairs; its only components are the sectional curvatures $K_{AB}=R^{AB}{}_{AB}$ (no sum), with $i,j\in\lbrace1,2,3\rbrace$, $k,l\in\lbrace4,\dots,7\rbrace$:
\begin{equation*}
K_{0i}=-(H_a'+H_a^2),\quad K_{0k}=-(H_b'+H_b^2),\quad K_{ij}=-H_a^2,\quad K_{kl}=-H_b^2,\quad K_{ik}=-H_aH_b .
\end{equation*}
The Einstein tensor and the scalar curvature are
\begin{align*}
G^0{}_0&=3H_a^2+12H_aH_b+6H_b^2, \\
G^1{}_1=G^2{}_2=G^3{}_3&=2H_a'+3H_a^2+4H_b'+10H_b^2+8H_aH_b, \\
G^4{}_4=\cdots=G^7{}_7&=3H_a'+6H_a^2+3H_b'+6H_b^2+9H_aH_b, \\
R&=-2\big(3H_a'+6H_a^2+4H_b'+10H_b^2+12H_aH_b\big).
\end{align*}
For the inhomogeneous metric the notebook checks, exactly at $x_0=1/2$, $x_4=1/3$, the symmetries of the Riemann tensor ($R_{\rho\sigma\mu\nu}=-R_{\sigma\rho\mu\nu}=-R_{\rho\sigma\nu\mu}=R_{\mu\nu\rho\sigma}$ with $R_{\rho\sigma\mu\nu}=g_{\rho\lambda}R^\lambda{}_{\sigma\mu\nu}$), the first Bianchi identity $R^\rho{}_{\sigma\mu\nu}+R^\rho{}_{\mu\nu\sigma}+R^\rho{}_{\nu\sigma\mu}=0$, and the contracted Bianchi identity
\begin{equation*}
\nabla_\mu G^{\mu}{}_{\nu}=\partial_\mu G^\mu{}_\nu+\{{}^{\mu}_{\mu\lambda}\}G^\lambda{}_\nu-\{{}^{\lambda}_{\mu\nu}\}G^\mu{}_\lambda=0 .
\end{equation*}
(the contracted Bianchi identity is also checked for the cosmological metric with arbitrary $a(x_0)$, $b(x_0)$). A fourth metric enters in section 5.4: the static metric $ds^2=f(r)\,dx_0^2-h(r)\,dr^2-r^2(dx_2^2+\cdots+dx_7^2)$ with $r=x_1$ and arbitrary positive $f$, $h$.

\Needspace{12\baselineskip}\section*{4. The Lovelock tensors of 7+1 spacetime (section 12)}\addcontentsline{toc}{section}{4. The Lovelock tensors of 7+1 spacetime (section 12)}
\Needspace{8\baselineskip}\subsection*{4.1 Lovelock’s theorem and the definitions}\addcontentsline{toc}{subsection}{4.1 Lovelock’s theorem and the definitions}
In $n$ dimensions the most general tensor $A^{lh}$ that is symmetric, divergence-free and built from the metric and its first and second derivatives is (D. Lovelock and H. Rund, \emph{Tensors, Differential Forms, and Variational Principles}, Dover; Remark 2 after their Eq. (4.38))
\begin{equation*}
A^{lh}=\sqrt{|g|}\sum_{k=1}^{m-1}\alpha_{(k)}\,g^{jl}\,\delta^{h h_1\cdots h_{2k}}_{j j_1\cdots j_{2k}}\,R^{j_1j_2}{}_{h_1h_2}\cdots R^{j_{2k-1}j_{2k}}{}_{h_{2k-1}h_{2k}}+\lambda\sqrt{|g|}\,g^{lh},\qquad m=\begin{cases}n/2,&n\ \text{even},\\(n+1)/2,&n\ \text{odd},\end{cases}
\end{equation*}
with arbitrary constants $\alpha_{(k)}$ and $\lambda$. For $n=8$, $m-1=8/2-1=3$: there are exactly three curvature terms. The generalized Kronecker delta is the determinant of the matrix of ordinary Kronecker deltas,
\begin{equation*}
\delta^{\mu_1\cdots\mu_n}_{\nu_1\cdots\nu_n}=\det\big[\delta^{\mu_i}_{\nu_j}\big]_{i,j=1}^{n},\qquad \delta^{\mu_1\cdots\mu_p\rho}_{\nu_1\cdots\nu_p\rho}=(8-p)\,\delta^{\mu_1\cdots\mu_p}_{\nu_1\cdots\nu_p}\ \ (D=8),
\end{equation*}
it is $+1$ ($-1$) when the lower indices are an even (odd) permutation of the upper ones, all different, and $0$ otherwise; it vanishes when $n>8$. With the normalization of the reference notebook Einstein-Lovelock-7+1-Nash.nb (P. L. Nash, GPL), which this document follows, the Lovelock tensors (\texttt{Lovelock1}, \texttt{Lovelock2}, \texttt{Lovelock3}) and densities are
\begin{align*}
\mathrm{Lovelock}_k\equiv E^{(k)\mu}{}_{\nu}&=-\frac{1}{2^{k+1}}\,\delta^{\mu\alpha_1\beta_1\cdots\alpha_k\beta_k}_{\nu\gamma_1\delta_1\cdots\gamma_k\delta_k}\,R^{\gamma_1\delta_1}{}_{\alpha_1\beta_1}\cdots R^{\gamma_k\delta_k}{}_{\alpha_k\beta_k}, \\
L_k&=\frac{1}{2^k}\,\delta^{\alpha_1\beta_1\cdots\alpha_k\beta_k}_{\gamma_1\delta_1\cdots\gamma_k\delta_k}\,R^{\gamma_1\delta_1}{}_{\alpha_1\beta_1}\cdots R^{\gamma_k\delta_k}{}_{\alpha_k\beta_k}.
\end{align*}
Special cases: $L_0=1$, $E^{(0)\mu}{}_\nu=-\tfrac12\delta^\mu_\nu$; $L_1=R$, $E^{(1)}=G$; $L_2=R^2-4R_{\mu\nu}R^{\mu\nu}+R_{\mu\nu\rho\sigma}R^{\mu\nu\rho\sigma}$ (Gauss–Bonnet) with $E^{(2)}$ the Lanczos tensor
\begin{equation*}
E^{(2)}_{\mu\nu}=2\big(R\,R_{\mu\nu}-2R_{\mu\alpha}R^{\alpha}{}_{\nu}-2R^{\alpha\beta}R_{\mu\alpha\nu\beta}+R_{\mu}{}^{\alpha\beta\gamma}R_{\nu\alpha\beta\gamma}\big)-\tfrac12 g_{\mu\nu}L_2 ;
\end{equation*}
$L_3$, $E^{(3)}$ are cubic in the curvature (section 4.3 gives $E^{(3)}$ explicitly). Because $E^{(k)}$ contains $2k+1$ different upper indices, $E^{(4)}\equiv0$ in eight dimensions, although $L_4$ (the Euler density) is not zero; $\sqrt{|g|}L_4$ is a total derivative. Every $E^{(k)}$ is symmetric with both indices down, divergence-free for every metric, and has the trace
\begin{equation*}
E^{(k)\mu}{}_{\mu}=-\frac{8-2k}{2}\,L_k .
\end{equation*}
\Needspace{8\baselineskip}\subsection*{4.2 Action and field equations}\addcontentsline{toc}{subsection}{4.2 Action and field equations}
The reference notebook writes the Einstein–Lovelock VACUUM field equations as
\begin{equation*}
0=-\Lambda\,\delta^\mu_\nu+H^{-2}\,w_1\,\mathrm{Lovelock}_1+H^{-4}\,w_2\,\mathrm{Lovelock}_2+H^{-6}\,w_3\,\mathrm{Lovelock}_3,
\end{equation*}
with pure numbers $w_1,w_2,w_3,\Lambda$ and a fundamental inverse length $H$ (it exists because the Lovelock tensors of different orders have different dimensions, length$^{-2k}$). With matter, and with $\delta(\sqrt{|g|}L_k)=\sqrt{|g|}\,E^{(k)}_{\mu\nu}\,\delta g^{\mu\nu}$ (up to a total derivative) and $\delta S_{\text{matter}}=\tfrac12\int\sqrt{|g|}\,T_{\mu\nu}\,\delta g^{\mu\nu}\,d^8X$, the action
\begin{equation*}
S=-\frac{1}{2\kappa}\int d^8X\,\sqrt{|g|}\,\sum_{k=0}^{3}\alpha_k L_k+S_{\text{matter}},\qquad \alpha_0=2\bar\Lambda,\ \ \alpha_1=1,
\end{equation*}
gives the Einstein–Lovelock field equations
\begin{equation*}
\sum_{k=0}^{3}\alpha_k\,E^{(k)\mu}{}_{\nu}=\kappa\,T^\mu{}_\nu\qquad\Longleftrightarrow\qquad G^\mu{}_\nu+\alpha_2E^{(2)\mu}{}_\nu+\alpha_3E^{(3)\mu}{}_\nu-\bar\Lambda\,\delta^\mu_\nu=\kappa\,T^\mu{}_\nu,
\end{equation*}
with $\kappa=8\pi G_8$ ($G_8$ the Newton constant of eight dimensions), the Gauss–Bonnet and cubic couplings $\alpha_2$, $\alpha_3$ (dimensions length² and length\texorpdfstring{\ensuremath{{}^{4}}}{⁴}) and the cosmological constant $\bar\Lambda$ (dimension length$^{-2}$; it is written $\Lambda$ in sections 7 and following, where no pure number $\Lambda$ appears). $\alpha_4$ does not appear: its tensor vanishes. For $T=0$ the two forms are the same equation (divide the first by $w_1/H^2$):
\begin{equation*}
\alpha_2=\frac{w_2}{w_1H^2},\qquad \alpha_3=\frac{w_3}{w_1H^4},\qquad \bar\Lambda=\frac{\Lambda H^2}{w_1}.
\end{equation*}
\Needspace{8\baselineskip}\subsection*{4.3 The Lovelock tensors with xAct, as in the reference notebook}\addcontentsline{toc}{subsection}{4.3 The Lovelock tensors with xAct, as in the reference notebook}
Section 12 of 7+1.nb computes the Lovelock tensors exactly as the reference notebook does, with the tensor packages xAct (xPerm, xTensor, xCoba, xTras; J. M. Martín-García et al., GPL; versions xPerm 1.2.3, xTensor 1.2.0, xCoba 0.8.6, xTras 1.4.2). The reference notebook also contains a "Model metric" which its author has marked FAILED, DO NOT USE (it never worked for the 7+1 signature); it is not used anywhere in this project. The set-up (abstract indices $a$, $b$, $c$, $d$, $f$, $i$, $j$, $k$, $l$, $m$, $n$, $p$, $q$, $r$, $s$, $y$; a metric \texttt{g71} with Levi-Civita derivative \texttt{CD}; a chart \texttt{chartM8} with the coordinate scalars \texttt{x0[]}, …, \texttt{x7[]}):

\begin{Verbatim}[fontsize=\small,frame=single,framesep=2mm]
Needs["xAct`xCoba`"]; Needs["xAct`xTras`"];
$DefInfoQ = False; $CVVerbose = False;
SetOptions[ToCanonical, UseMetricOnVBundle -> None, Verbose -> False];
DefManifold[M8, 8, {a, b, c, d, f, i, j, k, l, m, n, p, q, r, s, y}];
DefMetric[{1, 7, 0}, g71[-a, -b], CD, {";", "D"}, WeightedWithBasis -> AIndex];
DefChart[chartM8, M8, {0, 1, 2, 3, 4, 5, 6, 7}, {x0[], x1[], x2[], x3[], x4[], x5[], x6[],
    ...x7[]},
  ExtendedCoordinateDerivatives -> True, FormatBasis -> {"Partials", "Differentials"}];
MyArrayComponents[expr_] := expr // ToBasis[chartM8] // ComponentArray // ToValues // ToValues
    ...// Simplify;
kδ[lower_, upper_] /; Length[lower] == Length[upper] := Det[Outer[delta, lower, upper]];
\end{Verbatim}
In the xAct notation the signature is $\{p,q,z\}$ = (number of $+1$, of $-1$, of $0$ eigenvalues). The reference notebook declares $\{7,1,0\}$ (the mostly-plus convention); this notebook declares $\{1,7,0\}$, which matches its metric $\eta=\mathrm{diag}(+1,-1,\dots,-1)$. The declaration does not enter the Lovelock tensors with mixed indices, which follow from the components alone, and both declarations give xAct the same sign of the determinant of the metric, $(-1)^q=-1$ for $q=7$ and for $q=1$ (\texttt{SignDetOfMetric}). While the packages load, their banners and the warnings that xTensor cannot attach rules to \texttt{Commutator} and \texttt{Anticommutator} (protected built-in symbols of Wolfram 15, not used here) are suppressed. \texttt{kδ[lower, upper]} is the generalized Kronecker delta $\delta^{\text{upper}}_{\text{lower}}$ with abstract indices; the Lovelock tensors are then typed exactly as in the reference notebook,

\begin{Verbatim}[fontsize=\small,frame=single,framesep=2mm]
Lovelock1 = -(1/4) RiemannCD[-c,-d,f,i] kδ[{-a,-f,-i},{b,c,d}] // Simplification
Lovelock2 = -(1/8) RiemannCD[-c,-d,f,i] RiemannCD[-j,-k,l,m] kδ[{-a,-f,-i,-l,-m},{b,c,d,j,k}]
    ...// Simplification
Lovelock3 = -(1/16) RiemannCD[-c,-d,f,i] RiemannCD[-j,-k,l,m] RiemannCD[-n,-p,q,r]
              kδ[{-a,-f,-i,-l,-m,-q,-r},{b,c,d,j,k,n,p}] // Simplification
\end{Verbatim}
i.e. $\mathrm{Lovelock}_k=E^{(k)b}{}_a$ with the free indices $a$ (down) and $b$ (up) and $R_{cd}{}^{fi}$ = \texttt{RiemannCD[-c,-d,f,i]}. Simplification (xTras) reduces them to curvature monomials; with $R_{ab}{}^{cd}$ the Riemann tensor, $R_a{}^b$ the Ricci tensor and $R$ the scalar curvature the notebook obtains
\begin{align*}
\mathrm{Lovelock}_{1}&=- \frac{\delta_{a}{}^{b} R}{2} + R_{a}{}^{b},
\end{align*}
\begin{align*}
\mathrm{Lovelock}_{2}&=- \frac{\delta_{a}{}^{b} R^{2}}{2} + 2 \delta_{a}{}^{b} R^{j}{}_{k} R_{j}{}^{k} \\
&\quad - \frac{\delta_{a}{}^{b} R_{jk}{}^{lm} R^{jk}{}_{lm}}{2} + 2 R R_{a}{}^{b} \\
&\quad - 4 R_{a}{}^{j} R^{b}{}_{j} - 4 R^{j}{}_{k} R_{aj}{}^{bk} \\
&\quad + 2 R_{aj}{}^{kl} R^{bj}{}_{kl},
\end{align*}
\begin{align*}
\mathrm{Lovelock}_{3}&=- \frac{\delta_{a}{}^{b} R^{3}}{2} + 6 \delta_{a}{}^{b} R R^{j}{}_{k} R_{j}{}^{k} \\
&\quad - \frac{3 \delta_{a}{}^{b} R R_{jk}{}^{lm} R^{jk}{}_{lm}}{2} \\
&\quad - 8 \delta_{a}{}^{b} R_{j}{}^{l} R^{j}{}_{k} R^{k}{}_{l} \\
&\quad - 12 \delta_{a}{}^{b} R^{j}{}_{k} R^{l}{}_{m} R_{jl}{}^{km} \\
&\quad + 12 \delta_{a}{}^{b} R^{j}{}_{k} R_{jl}{}^{mn} R^{kl}{}_{mn} \\
&\quad + 4 \delta_{a}{}^{b} R^{jk}{}_{lm} R_{jn}{}^{lp} R_{kp}{}^{mn} \\
&\quad - \delta_{a}{}^{b} R^{jk}{}_{lm} R_{jk}{}^{np} R^{lm}{}_{np} + 3 R^{2} R_{a}{}^{b} \\
&\quad - 12 R R_{a}{}^{j} R^{b}{}_{j} - 12 R R^{j}{}_{k} R_{aj}{}^{bk} \\
&\quad + 6 R R_{aj}{}^{kl} R^{bj}{}_{kl} + 24 R_{j}{}^{l} R^{j}{}_{k} R_{al}{}^{bk} \\
&\quad + 24 R_{a}{}^{j} R^{b}{}_{k} R_{j}{}^{k} + 24 R_{a}{}^{j} R^{k}{}_{l} R^{bl}{}_{jk} \\
&\quad - 12 R_{a}{}^{j} R^{bk}{}_{lm} R_{jk}{}^{lm} - 12 R^{j}{}_{k} R_{a}{}^{b} R_{j}{}^{k} \\
&\quad - 12 R^{j}{}_{k} R_{aj}{}^{lm} R^{bk}{}_{lm} \\
&\quad - 24 R^{j}{}_{k} R_{al}{}^{km} R^{bl}{}_{jm} \\
&\quad + 24 R^{j}{}_{k} R_{al}{}^{bm} R_{jm}{}^{kl} \\
&\quad + 3 R_{a}{}^{b} R_{jk}{}^{lm} R^{jk}{}_{lm} + 24 R^{k}{}_{l} R^{b}{}_{j} R_{ak}{}^{jl} \\
&\quad - 12 R^{b}{}_{j} R_{ak}{}^{lm} R^{jk}{}_{lm} \\
&\quad - 24 R_{aj}{}^{kl} R^{bm}{}_{kn} R^{jn}{}_{lm} \\
&\quad + 6 R_{aj}{}^{kl} R^{bj}{}_{mn} R_{kl}{}^{mn} \\
&\quad - 12 R_{kl}{}^{mn} R^{jl}{}_{mn} R_{aj}{}^{bk}.
\end{align*}
Two checks are made with abstract indices (ToCanonical): $\mathrm{Lovelock}_1-(R_a{}^b-\tfrac12\delta_a{}^bR)=0$, and the trace $\mathrm{Lovelock}_2{}_a{}^a=-2L_2$ with $L_2=R^2-4R_{cd}R^{cd}+R_{cdfi}R^{cdfi}$. The reference notebook sets the option \texttt{UseMetricOnVBundle -> None} of ToCanonical (indices are not moved up and down); for the trace identity the metric must identify the dummy pairs $R_{cd}R^{cd}$ and $R^c{}_dR_c{}^d$, so that one call uses \texttt{UseMetricOnVBundle -> All}.

\Needspace{8\baselineskip}\subsection*{4.4 Components with xAct: CD = LC[g]}\addcontentsline{toc}{subsection}{4.4 Components with xAct: CD = LC[g]}
For a metric given as an 8×8 matrix of components with lower indices the reference notebook makes it a CTensor in the basis of \texttt{chartM8}, lets MetricCompute compute and simplify its curvature, and sets \texttt{CD = LC[g]} ("xAct magic": the abstract covariant derivative becomes the Levi-Civita connection of the component metric). After \texttt{RiemannCD = Riemann[CD]}, \texttt{RicciCD = Ricci[CD]}, \texttt{RicciScalarCD = RicciScalar[CD]} and \texttt{EinsteinCD = Einstein[CD]} every abstract expression evaluates to components. The notebook function \texttt{xActLovelock[gm]} does exactly these steps:

\begin{Verbatim}[fontsize=\small,frame=single,framesep=2mm]
xActLovelock[gm_] := Module[{res},
   g71 = CTensor[gm /. toXc, {-chartM8, -chartM8}];
   MetricCompute[g71, chartM8, All, Parallelize -> False, Verbose -> False];
   CD = LC[g71];
   RiemannCD = Riemann[CD]; RicciCD = Ricci[CD]; RicciScalarCD = RicciScalar[CD]; EinsteinCD =
    ...Einstein[CD];
   res = Transpose[asMat8[MyArrayComponents[#]]] & /@ {Lovelock1, Lovelock2, Lovelock3,
    ...EinsteinCD[-a, b]};
   res /. fromXc]
\end{Verbatim}
Here \texttt{toXc} replaces $x_\mu$ by the coordinate scalar \texttt{xμ[]} and \texttt{fromXc} undoes it, \texttt{asMat8} turns a vanishing result into the 8×8 zero matrix, and the Transpose makes row $\mu$, column $\nu$ hold $E^{(k)\mu}{}_\nu$ (the array of MyArrayComponents has the lower index $a$ first; the pp-wave of section 4.7, whose $E^{(1)\mu}{}_\nu$ is not a symmetric matrix, fixes this order). For the cosmological metric \texttt{xActLovelock} gives $\mathrm{Lovelock}_1=$ \texttt{EinsteinCD} $=$ the Einstein tensor of section 3, and the closed forms of section 4.6.

\Needspace{8\baselineskip}\subsection*{4.5 An independent computation: the pair-set algorithm}\addcontentsline{toc}{subsection}{4.5 An independent computation: the pair-set algorithm}
The same tensors are also computed without xAct. A direct evaluation of the generalized delta would run over $8^{4k+2}$ index combinations; the notebook functions \texttt{pairSets}, \texttt{orderedSeqs}, \texttt{lovelockE}, \texttt{lovelockL} (and tools/lovelock71.py in Python) use the antisymmetries instead. Let $\mathcal U$ run over the SETS of $k$ disjoint index pairs $(\alpha<\beta)$ not containing $\mu$, and $\mathcal V$ over the ORDERED sequences of $k$ disjoint pairs $(\gamma<\delta)$ that exhaust $(\mathcal U\cup\lbrace\mu\rbrace)\setminus\lbrace\nu\rbrace$. Each pair contributes a factor 2 (both orders of a pair give the same term), the $k!$ orderings of the pairs of $\mathcal U$ give the same sum, and
\begin{equation*}
E^{(k)\mu}{}_{\nu}=-2^{k-1}k!\sum_{\mathcal U}\sum_{\mathcal V}\mathrm{sgn}(\mu\,\mathcal U\to\nu\,\mathcal V)\prod_{i=1}^{k}R^{\mathcal V_i}{}_{\mathcal U_i},\qquad L_k=2^kk!\sum_{\mathcal U}\sum_{\mathcal V}\mathrm{sgn}(\mathcal U\to\mathcal V)\prod_{i=1}^{k}R^{\mathcal V_i}{}_{\mathcal U_i},
\end{equation*}
where sgn is the sign of the permutation that turns the upper index sequence into the lower one (in Mathematica \texttt{Signature[up] Signature[dn]}). For $k=3$ and fixed $\mu,\nu$ there are 105 upper pair-sets and 90 ordered lower sequences. The notebook compares $E^{(1)}$ with the brute-force contraction $-\tfrac14\delta^{\mu\alpha\beta}_{\nu\gamma\delta}R^{\gamma\delta}{}_{\alpha\beta}$, $E^{(2)}$ with the Lanczos tensor, and $E^{(1)},E^{(2)},E^{(3)}$ of the cosmological metric and of the pp-wave with xAct: the two computations agree exactly. The pair-set sum also works on a table of exact rational numbers (section 4.7).

\Needspace{8\baselineskip}\subsection*{4.6 Closed forms for the cosmological metric}\addcontentsline{toc}{subsection}{4.6 Closed forms for the cosmological metric}
For the cosmological metric every $E^{(k)}$ is diagonal; the entry $E^{(k)\mu}{}_\mu$ (no sum) is $-\tfrac12$ times the density $L_k$ restricted to the index pairs that avoid $\mu$. The results (Ha = $H_a$, dHa = $H_a'$, and so on) are:
\begin{align*}
E^{(1)0}{}_0&=3 \left(H_{a}^{2} + 4 H_{a} H_{b} + 2 H_{b}^{2}\right), \\
E^{(1)1}{}_1&=2 H'_{a} + 4 H'_{b} + 3 H_{a}^{2} + 8 H_{a} H_{b} + 10 H_{b}^{2}, \\
E^{(1)4}{}_4&=3 H'_{a} + 3 H'_{b} + 6 H_{a}^{2} + 9 H_{a} H_{b} + 6 H_{b}^{2}.
\end{align*}
\begin{align*}
E^{(2)0}{}_0&=- 12 H_{b} \left(4 H_{a}^{3} + 18 H_{a}^{2} H_{b} + 12 H_{a} H_{b}^{2} + H_{b}^{3}\right), \\
E^{(2)1}{}_1&=- 32 H'_{a} H_{a} H_{b} - 48 H'_{a} H_{b}^{2} - 16 H'_{b} H_{a}^{2} - 96 H'_{b} H_{a} H_{b} \\
&\quad - 48 H'_{b} H_{b}^{2} - 32 H_{a}^{3} H_{b} - 136 H_{a}^{2} H_{b}^{2} - 192 H_{a} H_{b}^{3} \\
&\quad - 60 H_{b}^{4}, \\
E^{(2)4}{}_4&=- 12 H'_{a} H_{a}^{2} - 72 H'_{a} H_{a} H_{b} - 36 H'_{a} H_{b}^{2} - 36 H'_{b} H_{a}^{2} \\
&\quad - 72 H'_{b} H_{a} H_{b} - 12 H'_{b} H_{b}^{2} - 12 H_{a}^{4} - 108 H_{a}^{3} H_{b} \\
&\quad - 180 H_{a}^{2} H_{b}^{2} - 108 H_{a} H_{b}^{3} - 12 H_{b}^{4}.
\end{align*}
\begin{align*}
E^{(3)0}{}_0&=360 H_{a}^{2} H_{b}^{3} \left(4 H_{a} + 3 H_{b}\right), \\
E^{(3)1}{}_1&=72 H_{b}^{2} \left(8 H'_{a} H_{a} H_{b} + 2 H'_{a} H_{b}^{2} + 12 H'_{b} H_{a}^{2} + 8 H'_{b} H_{a} H_{b} + 8 H_{a}^{3} H_{b} + 19 H_{a}^{2} H_{b}^{2} + 8 H_{a} H_{b}^{3}\right), \\
E^{(3)4}{}_4&=72 H_{a} H_{b} \left(9 H'_{a} H_{a} H_{b} + 6 H'_{a} H_{b}^{2} + 6 H'_{b} H_{a}^{2} + 9 H'_{b} H_{a} H_{b} + 9 H_{a}^{3} H_{b} + 17 H_{a}^{2} H_{b}^{2} + 9 H_{a} H_{b}^{3}\right).
\end{align*}
In particular
\begin{equation*}
E^{(2)0}{}_0=-12H_b\big(4H_a^3+18H_a^2H_b+12H_aH_b^2+H_b^3\big),\qquad E^{(3)0}{}_0=360\,H_a^2H_b^3\big(4H_a+3H_b\big).
\end{equation*}
Divergence-freeness becomes, for each $k$,
\begin{equation*}
\frac{d}{dx_0}E^{(k)0}{}_0+3H_a\big(E^{(k)0}{}_0-E^{(k)1}{}_1\big)+4H_b\big(E^{(k)0}{}_0-E^{(k)4}{}_4\big)=0,
\end{equation*}
checked for arbitrary $a,b$. The VARIATIONAL ORIGIN is checked as well: with a lapse function $N(x_0)$ ($g_{00}=N^2$) the reduced Lagrangian is $N a^3b^4L_k$, and its Euler–Lagrange expressions at $N=1$ are
\begin{equation*}
\frac{\delta(Na^3b^4L_k)}{\delta N}=-2a^3b^4E^{(k)0}{}_0,\qquad \frac{\delta(Na^3b^4L_k)}{\delta a}=-6a^2b^4E^{(k)1}{}_1,\qquad \frac{\delta(Na^3b^4L_k)}{\delta b}=-8a^3b^3E^{(k)4}{}_4,
\end{equation*}
for $k=1,2,3,4$ (for $k=4$ all three vanish identically: $\sqrt{|g|}L_4$ is a total derivative). The factors come from $\delta g^{00}=-2\delta N$, three $\delta g^{ii}=2\delta a/a^3$ and four $\delta g^{kk}=2\delta b/b^3$.

\Needspace{8\baselineskip}\subsection*{4.7 An inhomogeneous metric and a gravitational wave}\addcontentsline{toc}{subsection}{4.7 An inhomogeneous metric and a gravitational wave}
For the explicit inhomogeneous metric $g=\mathrm{diag}(n^2,-a^2,-a^2,-a^2,-b^2,-b^2,-b^2,-b^2)$ with $n=1+x_0x_4/5+x_0^2/11$, $a=1+x_0/3+x_4^2/7$, $b=2-x_0x_4/4+x_0^2/9+x_4/13$, the divergence $\nabla_\mu E^{(k)\mu}{}_\nu$ is evaluated EXACTLY (rational arithmetic) at $x_0=1/2$, $x_4=1/3$ for $k=1,2,3$; it vanishes. The derivative of $E^{(k)}$ (a polynomial of degree $k$ in the curvature) is obtained from the curvature and its first derivatives at the point: $\partial_\lambda E^{(k)}=\frac{d}{d\epsilon}E^{(k)}[R+\epsilon\,\partial_\lambda R]\big|_{\epsilon=0}$. At that point $E^{(2)}$ equals the Lanczos tensor exactly, and
\begin{gather*}
E^{(3)0}{}_4\big|_{x_0=1/2,\,x_4=1/3}=q\approx 0.0143928243619,\\ q\in\mathbb Q\ \text{(an exact fraction: numerator of 50, denominator of 52 digits)} .
\end{gather*}
For the pp-wave $ds^2=dx_0^2-dx_1^2-\cdots-dx_7^2+F(x_0-x_1,x_2,x_3)\,(dx_0-dx_1)^2$ with an arbitrary profile $F$ and the null vector $l_\mu=(1,-1,0,\dots,0)$,
\begin{equation*}
E^{(1)}_{\mu\nu}=\tfrac12\big(\partial_2^2+\partial_3^2\big)F\;l_\mu l_\nu,\qquad E^{(2)}=E^{(3)}=0\ \text{for every profile }F,
\end{equation*}
so a plane wave whose profile is harmonic in the transverse directions solves the Einstein–Lovelock equations for every $\alpha_2,\alpha_3$ (with $\bar\Lambda=0$). xAct and the pair-set sum agree for this metric, which is not diagonal.

\Needspace{8\baselineskip}\subsection*{4.8 Exponential vacuum solutions: eternal inflation of ordinary space and deflation of the auxiliary space}\addcontentsline{toc}{subsection}{4.8 Exponential vacuum solutions: eternal inflation of ordinary space and deflation of the auxiliary space}
The procedure of the reference notebook — solve the vacuum equations of section 4.2 for $w_2$, $w_3$ (and $\Lambda$) — is applied to the cosmological metric with constant Hubble rates, $a=e^{H_a x_0}$, $b=e^{H_b x_0}$, $H_a=H$, $H_b=q_HH$ ($q_H$ a pure number; code \texttt{qH}). With $H_a'=H_b'=0$ in the closed forms of section 4.6 every $E^{(k)}$ is a constant diagonal matrix proportional to $H^{2k}$, and the 64 equations reduce to three distinct ones (components 00, 11, 44). They are not independent: for constant Hubble rates the divergence identity of section 4.6 becomes
\begin{equation*}
3H_a\big(\mathrm{eq}_{00}-\mathrm{eq}_{11}\big)+4H_b\big(\mathrm{eq}_{00}-\mathrm{eq}_{44}\big)=0,
\end{equation*}
so for every $(H_a,H_b)$ there is a one-parameter family of couplings: given $w_1$ and $\Lambda$ the equations fix $w_2$, $w_3$. For $q_H=-1$ — ordinary space expands as $e^{Hx_0}$ while the auxiliary space contracts as $e^{-Hx_0}$, for every time $x_0$ (eternal inflation of ordinary space with eternal deflation of the auxiliary space) — the three equations are
\begin{equation*}
\mathrm{eq}_{00}=- \Lambda - 3 w_{1} - 36 w_{2} - 360 w_{3},\qquad \mathrm{eq}_{11}=- \Lambda + 5 w_{1} + 28 w_{2} + 216 w_{3},\qquad \mathrm{eq}_{44}=- \Lambda + 3 w_{1} + 12 w_{2} + 72 w_{3},
\end{equation*}
and the family is
\begin{gather*}
w_2=\frac{\Lambda}{4} - \frac{w_{1}}{2},\qquad w_3=- \frac{\Lambda}{36} + \frac{w_{1}}{24};\\ \Lambda=0:\quad \frac{w_2}{w_1}=- \frac{1}{2},\ \ \frac{w_3}{w_1}=\frac{1}{24},\quad\text{i.e.}\quad \alpha_2=\frac{w_2}{w_1H^2}=- \frac{1}{2}\,H^{-2},\ \ \alpha_3=\frac{w_3}{w_1H^4}=\frac{1}{24}\,H^{-4}.
\end{gather*}
For general $q_H$ (computed here from the closed forms of section 4.6; the notebook and the Python program obtain the same expressions, and check that they solve all 64 components for every $\Lambda$):
\begin{equation*}
w_2=\frac{\big(12 q_{H}^{3} + 27 q_{H}^{2} + 22 q_{H} + 9\big)\,w_1-\big(2 q_{H} + 3\big)\,\Lambda}{4 q_{H} \left(6 q_{H}^{4} + 36 q_{H}^{3} + 66 q_{H}^{2} + 51 q_{H} + 16\right)},
\end{equation*}
\begin{equation*}
w_3=\frac{\big(15 q_{H}^{4} + 36 q_{H}^{3} + 33 q_{H}^{2} + 18 q_{H} + 3\big)\,w_1-\big(3 q_{H}^{2} + 6 q_{H} + 1\big)\,\Lambda}{72 q_{H}^{3} \left(6 q_{H}^{4} + 36 q_{H}^{3} + 66 q_{H}^{2} + 51 q_{H} + 16\right)}.
\end{equation*}
\Needspace{12\baselineskip}\section*{5. Generalized frame matrices, canonical spin connection, total covariant derivative (section 13)}\addcontentsline{toc}{section}{5. Generalized frame matrices, canonical spin connection, total covariant derivative (section 13)}
\Needspace{8\baselineskip}\subsection*{5.1 Frame field and generalized frame matrices}\addcontentsline{toc}{subsection}{5.1 Frame field and generalized frame matrices}
A frame field (vielbein) is a square root of the metric, and its inverse is written $e_A{}^\mu$:
\begin{equation*}
g_{\mu\nu}=\eta_{AB}\,e^A{}_\mu e^B{}_\nu,\qquad e_A{}^\mu e^A{}_\nu=\delta^\mu_\nu .
\end{equation*}
The GENERALIZED FRAME MATRICES are the position-dependent real 16×16 matrices
\begin{equation*}
\Gamma_\mu(X)=e^A{}_\mu(X)\,\Gamma_A,\qquad \Gamma^\mu(X)=e_A{}^\mu(X)\,\Gamma^A,\qquad \Gamma_\mu\Gamma_\nu+\Gamma_\nu\Gamma_\mu=2g_{\mu\nu}(X)\,1_{16},
\end{equation*}
with $\Gamma^A=\eta^{AA}\Gamma_A$ (no sum). They are real because the frame field and the $\Gamma_A$ are real.

\Needspace{8\baselineskip}\subsection*{5.2 The canonical spin connection and the total covariant derivative}\addcontentsline{toc}{subsection}{5.2 The canonical spin connection and the total covariant derivative}
The canonical (Levi-Civita) spin connection and its action on 16-component spinors are
\begin{align*}
\omega_\mu{}^A{}_B&=e^A{}_\nu\big(\partial_\mu e_B{}^\nu+\{{}^{\nu}_{\mu\lambda}\}e_B{}^\lambda\big),\qquad \omega_{\mu AB}=\eta_{AC}\,\omega_\mu{}^C{}_B=-\omega_{\mu BA}, \\
\Omega_\mu&=\tfrac14\,\omega_{\mu AB}\,\Gamma^A\Gamma^B,\qquad D_\mu\psi=\partial_\mu\psi+\Omega_\mu\psi,\qquad D_\mu\bar\psi=\partial_\mu\bar\psi-\bar\psi\,\Omega_\mu .
\end{align*}
The TOTAL covariant derivative acts on every index of an object: with the Christoffel symbols on coordinate indices, with $\omega$ on frame indices, and through the commutator with $\Omega$ on the two spinor indices of a 16×16 matrix. The canonical spin connection is the one for which the total covariant derivative of the frame field and of every frame matrix vanishes:
\begin{align*}
D_\mu e^A{}_\nu&=\partial_\mu e^A{}_\nu-\{{}^{\lambda}_{\mu\nu}\}e^A{}_\lambda+\omega_\mu{}^A{}_B\,e^B{}_\nu=0\qquad\text{(vielbein postulate)}, \\
\nabla_\mu\Gamma_\nu&=\partial_\mu\Gamma_\nu-\{{}^{\lambda}_{\mu\nu}\}\Gamma_\lambda+[\Omega_\mu,\Gamma_\nu]=0, \\
D_\mu\Gamma^A&=\omega_\mu{}^A{}_B\,\Gamma^B+[\Omega_\mu,\Gamma^A]=0 .
\end{align*}
$\Omega_\mu$ is real and $\Gamma_0\Omega_\mu^{\mathsf T}\Gamma_0=-\Omega_\mu$, which makes $D_\mu\bar\psi$ the Dirac adjoint of $D_\mu\psi$ and the action below real.

\Needspace{8\baselineskip}\subsection*{5.3 Spin curvature, Bianchi identity, Lichnerowicz}\addcontentsline{toc}{subsection}{5.3 Spin curvature, Bianchi identity, Lichnerowicz}
The curvature of the spin connection is the Riemann tensor in spinor form:
\begin{align*}
\mathcal F_{\mu\nu}&=\partial_\mu\Omega_\nu-\partial_\nu\Omega_\mu+[\Omega_\mu,\Omega_\nu]=[D_\mu,D_\nu]=\tfrac14\,R_{AB\mu\nu}\,\Gamma^A\Gamma^B,\qquad R_{AB\mu\nu}=\eta_{AC}\,e^C{}_\rho\,R^\rho{}_{\sigma\mu\nu}\,e_B{}^\sigma, \\
0&=\partial_\lambda \mathcal F_{\mu\nu}+[\Omega_\lambda,\mathcal F_{\mu\nu}]+\partial_\mu \mathcal F_{\nu\lambda}+[\Omega_\mu,\mathcal F_{\nu\lambda}]+\partial_\nu \mathcal F_{\lambda\mu}+[\Omega_\nu,\mathcal F_{\lambda\mu}]\qquad\text{(Bianchi)}, \\
\tfrac12\,\Gamma^\mu\Gamma^\nu \mathcal F_{\mu\nu}&=-\tfrac{R}{4}\,1_{16},\qquad (\Gamma^\mu D_\mu)^2\psi=g^{\mu\nu}\big(\partial_\mu D_\nu\psi+\Omega_\mu D_\nu\psi-\{{}^{\lambda}_{\mu\nu}\}D_\lambda\psi\big)-\tfrac{R}{4}\psi .
\end{align*}
Every identity of sections 5.1–5.3 is checked by the notebook (function \texttt{frameChecks}) and by the Python program for four metrics: the cosmological metric with the diagonal frame $e=\mathrm{diag}(1,a,a,a,b,b,b,b)$; the pp-wave with the non-diagonal frame $e^0=(1+F/2)dx_0-(F/2)dx_1$, $e^1=(F/2)dx_0+(1-F/2)dx_1$, $e^A=dx^A$ otherwise; the inhomogeneous metric with $e=\mathrm{diag}(n,a,a,a,b,b,b,b)$ (exactly at $x_0=1/2$, $x_4=1/3$, the derivatives being taken symbolically first); and the static radial metric of section 5.4.

\Needspace{8\baselineskip}\subsection*{5.4 Cartan’s first structure equation; a static metric with a radial coordinate}\addcontentsline{toc}{subsection}{5.4 Cartan’s first structure equation; a static metric with a radial coordinate}
With the frame one-forms $e^A=e^A{}_\mu dx^\mu$ and the connection one-forms $\omega^A{}_B=\omega_\mu{}^A{}_B\,dx^\mu$, the vielbein postulate antisymmetrized in $\mu,\nu$ (the Christoffel terms cancel) is Cartan’s first structure equation without torsion:
\begin{equation*}
de^A+\omega^A{}_B\wedge e^B=0\quad\Longleftrightarrow\quad \partial_\mu e^A{}_\nu-\partial_\nu e^A{}_\mu+\omega_\mu{}^A{}_B\,e^B{}_\nu-\omega_\nu{}^A{}_B\,e^B{}_\mu=0 .
\end{equation*}
It determines the canonical spin connection from the frame alone. With the structure coefficients $C^A{}_{BC}=-C^A{}_{CB}$ and the connection written with a frame form index,
\begin{align*}
de^A&=\tfrac12\,C^A{}_{BC}\,e^B\wedge e^C,\qquad C^A{}_{BC}=e_B{}^\mu e_C{}^\nu\big(\partial_\mu e^A{}_\nu-\partial_\nu e^A{}_\mu\big),\qquad C_{ABC}=\eta_{AD}C^D{}_{BC}, \\
\omega_{C,AB}&\equiv e_C{}^\mu\,\omega_{\mu AB}=-\tfrac12\big(C_{ACB}+C_{BAC}-C_{CBA}\big),
\end{align*}
the unique solution with $\omega_{C,AB}=-\omega_{C,BA}$; both relations are checked for all four frames. As an example take $r=x_1>0$ and arbitrary positive functions $f(r)$, $h(r)$ (code \texttt{fR[x1]}, \texttt{hR[x1]}):
\begin{gather*}
ds^2=f(r)\,dx_0^2-h(r)\,dr^2-r^2\big(dx_2^2+dx_3^2+dx_4^2+dx_5^2+dx_6^2+dx_7^2\big),\\ e^0=\sqrt f\,dx_0,\ \ e^1=\sqrt h\,dr,\ \ e^k=r\,dx^k\ (k=2,\dots,7).
\end{gather*}
The six directions $x_2,\dots,x_7$ are flat directions scaled by $r$ (a static "planar" metric of 7+1 spacetime, not a spherically symmetric one). Cartan’s equation gives $de^0=\frac{f'}{2\sqrt f}\,dr\wedge dx_0=-\omega^0{}_1\wedge e^1$ and $de^k=dr\wedge dx^k=-\omega^k{}_1\wedge e^1$, hence $\omega^0{}_1=\frac{f'}{2\sqrt{fh}}\,dx_0$, $\omega^k{}_1=\frac{1}{\sqrt h}\,dx^k$, and
\begin{align*}
\Omega_0&=\frac{f'}{4\sqrt f\sqrt h}\,\Gamma_1\Gamma_0,\qquad \Omega_1=0,\qquad \Omega_k=\frac{1}{2\sqrt h}\,\Gamma_1\Gamma_k\quad(k=2,\dots,7), \\
\Gamma^\mu\Omega_\mu&=\frac{1}{\sqrt h}\Big(\frac{f'}{4f}+\frac3r\Big)\Gamma^1,\qquad \Gamma^1=\eta^{11}\Gamma_1=-\Gamma_1,
\end{align*}
the $3/r=6\times\frac{1}{2r}$ coming from the six scaled directions. The covariant Dirac operator and a rescaling that removes the connection are therefore
\begin{align*}
i\Gamma^\mu D_\mu-m&=\frac{i}{\sqrt f}\,\Gamma^0\partial_0+\frac{i}{\sqrt h}\,\Gamma^1\Big(\partial_r+\frac{f'}{4f}+\frac3r\Big)+\sum_{k=2}^{7}\frac{i}{r}\,\Gamma^k\partial_k-m, \\
\psi&=f^{-1/4}r^{-3}\chi:\qquad \Gamma^\mu D_\mu\psi=f^{-1/4}r^{-3}\,\Gamma^\mu\partial_\mu\chi .
\end{align*}
\Needspace{12\baselineskip}\section*{6. The covariant Dirac equation of dirac16complex}\addcontentsline{toc}{section}{6. The covariant Dirac equation of dirac16complex}
The Lagrangian density and the field equation are
\begin{equation*}
\mathcal L=\sqrt{|g|}\;\bar\psi\big(i\,\Gamma^\mu D_\mu-m\big)\psi,\qquad \big(i\,\Gamma^\mu(X)\,D_\mu-m\big)\psi=0,
\end{equation*}
where the spin connection $\Omega_\mu$ in $D_\mu$ is the coupling of gravity to dirac16complex. The current $J^\mu=\bar\psi\Gamma^\mu\psi$ is conserved because, for arbitrary $\psi$ and $\bar\psi$ (by $\nabla\Gamma=0$ and $\partial_\mu\sqrt{|g|}=\sqrt{|g|}\{{}^{\lambda}_{\lambda\mu}\}$),
\begin{equation*}
\partial_\mu\big(\sqrt{|g|}\,\bar\psi\Gamma^\mu\psi\big)=\sqrt{|g|}\,\big[(D_\mu\bar\psi)\Gamma^\mu\psi+\bar\psi\Gamma^\mu D_\mu\psi\big],
\end{equation*}
and the bracket vanishes when $\psi$ and $\bar\psi$ solve the field equation and its adjoint. Squaring with the Lichnerowicz formula gives $\big(g^{\mu\nu}\nabla_\mu D_\nu+m^2-\tfrac{R}{4}\big)\psi=0$.

\textbf{Coupling to a gauge field.} An abelian gauge field $A_\mu(X)$ with charge $q$ enters through the gauge-covariant derivative
\begin{equation*}
D_\mu\psi=\partial_\mu\psi+\Omega_\mu\psi-iqA_\mu\psi,\qquad D_\mu\bar\psi=\partial_\mu\bar\psi-\bar\psi\,\Omega_\mu+iqA_\mu\bar\psi .
\end{equation*}
Under $\psi\to e^{iq\alpha}\psi$, $A_\mu\to A_\mu+\partial_\mu\alpha$ ($\alpha(X)$ arbitrary) the Dirac operator is covariant, $\Gamma^\mu D'_\mu(e^{iq\alpha}\psi)=e^{iq\alpha}\,\Gamma^\mu D_\mu\psi$, and the current identity above holds unchanged, so the electric current $q\,\bar\psi\Gamma^\mu\psi$ is conserved. Both statements are checked with the cosmological frame and arbitrary $A_\mu(x_0,x_1,x_5)$, $\alpha(x_0,x_1,x_5)$.

\Needspace{8\baselineskip}\subsection*{6.1 The cosmological spacetime}\addcontentsline{toc}{subsection}{6.1 The cosmological spacetime}
For the cosmological metric the spin connection is NOT zero whenever space expands:
\begin{equation*}
\Omega_0=0,\qquad \Omega_i=\tfrac{a'}{2}\,\Gamma_i\Gamma_0\ (i=1,2,3),\qquad \Omega_k=\tfrac{b'}{2}\,\Gamma_k\Gamma_0\ (k=4,\dots,7),\qquad \Gamma^\mu\Omega_\mu=\Big(\tfrac32H_a+2H_b\Big)\Gamma_0 .
\end{equation*}
The rescaled field $\chi=(a^3b^4)^{1/2}\psi$ obeys an equation without connection, $\Gamma^\mu D_\mu\psi=(a^3b^4)^{-1/2}\Gamma^\mu\partial_\mu\chi$. For a plane wave $\chi=e^{i(\vec p\cdot\vec x+\vec q\cdot\vec y)}c(x_0)$ — $\vec p=(p_1,p_2,p_3)$ the momentum along ordinary space $\vec x=(x_1,x_2,x_3)$, $\vec q=(q_4,\dots,q_7)$ along the auxiliary space $\vec y=(x_4,\dots,x_7)$ — the Dirac equation becomes
\begin{equation*}
i\,\frac{dc}{dx_0}=h(x_0)\,c,\qquad h(x_0)=\sum_{i=1}^{3}\alpha_i\,\frac{p_i}{a}+\sum_{k=4}^{7}\alpha_k\,\frac{q_k}{b}+\beta\,m,\qquad \beta=\Gamma_0,\ \ \alpha_j=\Gamma^0\Gamma^j=\begin{pmatrix}0&\tau^O_j\\-\tau^O_j&0\end{pmatrix},
\end{equation*}
the flat-space Hamiltonian with the physical (redshifted) momenta $\vec p/a$ and $\vec q/b$.

\Needspace{8\baselineskip}\subsection*{6.2 Conformal weight 7/2}\addcontentsline{toc}{subsection}{6.2 Conformal weight 7/2}
For a conformally flat metric $g=C(x_0)^2\eta$ ($x_0$ is then conformal time) the massless operator satisfies, for every $f$,
\begin{equation*}
\Gamma^\mu D_\mu\big(C^{-7/2}f\big)=C^{-9/2}\,\Gamma^A\partial_Af,
\end{equation*}
so a massless dirac16complex in a conformally flat 7+1 spacetime is a rescaled flat-space solution and no pairs are created. The cosmological metric is conformally flat only when $a=b$.

\Needspace{8\baselineskip}\subsection*{6.3 Two classes of states in an anisotropic expansion}\addcontentsline{toc}{subsection}{6.3 Two classes of states in an anisotropic expansion}
With unit vectors $\vec n_a$, $\vec n_b$ along $\vec p$ and $\vec q$ the Hamiltonian is $h=\frac{|\vec p|}{a}\mathcal A+\frac{|\vec q|}{b}\mathcal B+m\beta$ with the CONSTANT real matrices $\mathcal A=\sum_i n_{a,i}\alpha_i$, $\mathcal B=\sum_k n_{b,k}\alpha_k$. $\mathcal A$, $\mathcal B$, $\beta$ square to $1_{16}$ and anticommute; $\mathcal Z=\mathcal A\mathcal B\beta$ squares to $-1_{16}$, is antisymmetric, and commutes with $\mathcal A$, $\mathcal B$, $\beta$ and $h$, so its eigenvalue $\pm i$ is conserved. On $\mathcal Z=+i$ the three matrices act like the Pauli matrices $(\sigma_1,\sigma_2,\sigma_3)$, on $\mathcal Z=-i$ like $(\sigma_1,\sigma_2,-\sigma_3)$; the 16-component mode equation is four copies of each of
\begin{equation*}
i\,\frac{dc}{dx_0}=\Big(\frac{|\vec p|}{a}\sigma_1+\frac{|\vec q|}{b}\sigma_2\pm m\,\sigma_3\Big)c,\qquad c\in\mathbb C^2 .
\end{equation*}
Pairs are created when the direction of the physical momentum $(|\vec p|/a,|\vec q|/b)$ turns fast compared with the frequency $\omega=\sqrt{|\vec p|^2/a^2+|\vec q|^2/b^2+m^2}$. The Bogoliubov coefficient $|\beta|^2$ of a mode is the probability that a positive-energy state at the initial time ends in a negative-energy state at the final time; $|\alpha|^2+|\beta|^2=1$.

\Needspace{12\baselineskip}\section*{7. The three epochs (section 14)}\addcontentsline{toc}{section}{7. The three epochs (section 14)}
For the cosmological metric with $T^\mu{}_\nu=\mathrm{diag}(\varepsilon,-p_a,-p_a,-p_a,-p_b,-p_b,-p_b,-p_b)$ (energy density $\varepsilon$, pressures $p_a$ along ordinary space and $p_b$ along the auxiliary space) the field equations are
\begin{equation*}
\sum_k\alpha_kE^{(k)0}{}_0=\kappa\varepsilon,\qquad \sum_k\alpha_kE^{(k)1}{}_1=-\kappa p_a,\qquad \sum_k\alpha_kE^{(k)4}{}_4=-\kappa p_b,
\end{equation*}
and they imply the conservation law $\varepsilon'+3H_a(\varepsilon+p_a)+4H_b(\varepsilon+p_b)=0$.

\Needspace{8\baselineskip}\subsection*{7.1 Primordial epoch (near \texorpdfstring{$x_0=0$}{x0=0}): Kasner solutions and dynamical compactification}\addcontentsline{toc}{subsection}{7.1 Primordial epoch (near \texorpdfstring{$x_0=0$}{x0=0}): Kasner solutions and dynamical compactification}
For $a=x_0^{p}$, $b=x_0^{q}$ every term of $E^{(k)}$ scales as $x_0^{-2k}$: near $x_0=0$ the highest Lovelock order with a non-zero coupling dominates, and matter whose density grows more slowly than $x_0^{-2k}$ is negligible. With $H_a=p/x_0$, $H_b=q/x_0$, $H_a'=-p/x_0^2$, $H_b'=-q/x_0^2$ in the closed forms of section 4.6 the vacuum solutions of each order are:

\begin{itemize}
\item \textbf{Einstein} ($k=1$): $3p+4q=1$ and $3p^2+4q^2=1$, i.e. $p=(1\pm2\sqrt2)/7$, $q=(2\mp3\sqrt2)/14$. On the branch $p=0.5469$, $q=-0.1602$ ordinary space expands while the auxiliary space contracts as $x_0$ grows: a dynamical compactification of $x_4,\dots,x_7$.
\item \textbf{Gauss–Bonnet} ($k=2$): $3p+4q=3$, and $p$ is a root of $-203p^3+81p^2+351p+27=0$: $p=-1.0835018$, $-0.078631058$, $1.5611476$ (with $q=(3-3p)/4$).
\item \textbf{Cubic Lovelock} ($k=3$): the non-degenerate solution $(p,q)=(-15/7,\,20/7)$ with $3p+4q=5$; $E^{(3)}$ vanishes identically on the degenerate families $q=0$ (static auxiliary space) and $p=0$.
\end{itemize}
The pure-order equations are (each multiplied by $x_0^{2k}$; columns 00, 11, 44):
\begin{align*}
x_0^{2}E^{(1)0}{}_0&=3 \left(p^{2} + 4 p q + 2 q^{2}\right), \\
x_0^{2}E^{(1)1}{}_1&=3 p^{2} + 8 p q - 2 p + 10 q^{2} - 4 q, \\
x_0^{2}E^{(1)4}{}_4&=3 \left(2 p^{2} + 3 p q - p + 2 q^{2} - q\right).
\end{align*}
\begin{align*}
x_0^{4}E^{(2)0}{}_0&=- 12 q \left(4 p^{3} + 18 p^{2} q + 12 p q^{2} + q^{3}\right), \\
x_0^{4}E^{(2)1}{}_1&=- 4 q \left(8 p^{3} + 34 p^{2} q - 12 p^{2} + 48 p q^{2} - 36 p q + 15 q^{3} - 12 q^{2}\right), \\
x_0^{4}E^{(2)4}{}_4&=- 12 \left(p^{4} + 9 p^{3} q - p^{3} + 15 p^{2} q^{2} - 9 p^{2} q + 9 p q^{3} - 9 p q^{2} + q^{4} - q^{3}\right).
\end{align*}
\begin{align*}
x_0^{6}E^{(3)0}{}_0&=360 p^{2} q^{3} \left(4 p + 3 q\right), \\
x_0^{6}E^{(3)1}{}_1&=72 p q^{3} \left(8 p^{2} + 19 p q - 20 p + 8 q^{2} - 10 q\right), \\
x_0^{6}E^{(3)4}{}_4&=72 p^{2} q^{2} \left(9 p^{2} + 17 p q - 15 p + 9 q^{2} - 15 q\right).
\end{align*}
\Needspace{8\baselineskip}\subsection*{7.2 Pair creation in the primordial epoch}\addcontentsline{toc}{subsection}{7.2 Pair creation in the primordial epoch}
On the Einstein–Kasner branch a massless mode depends on $(|\vec p|,|\vec q|)$ only through the scaling variable $K=|\vec p|\,\lambda^{1-p}$, $\lambda=(|\vec q|/|\vec p|)^{1/(q-p)}$ ($x_0\to\lambda x_0$ maps one problem onto the other). The evolution is adiabatic before and after the time at which $|\vec p|/a\approx|\vec q|/b$, so it is integrated from $10^{-15}\lambda$ to $300\lambda$:

{\small
\begin{longtable}{>{\raggedright\arraybackslash}p{\dimexpr 0.2000\linewidth-2\tabcolsep\relax}>{\raggedright\arraybackslash}p{\dimexpr 0.2000\linewidth-2\tabcolsep\relax}>{\raggedright\arraybackslash}p{\dimexpr 0.3000\linewidth-2\tabcolsep\relax}>{\raggedright\arraybackslash}p{\dimexpr 0.3000\linewidth-2\tabcolsep\relax}}
\toprule
\textbf{$(|\vec p|,|\vec q|)$} & \textbf{$K$} & \textbf{$|\beta|^2$ (Mathematica)} & \textbf{$|\beta|^2$ (Python)} \\
\midrule\endfirsthead
\toprule
\textbf{$(|\vec p|,|\vec q|)$} & \textbf{$K$} & \textbf{$|\beta|^2$ (Mathematica)} & \textbf{$|\beta|^2$ (Python)} \\
\midrule\endhead
\endfoot
\bottomrule\endlastfoot
$(1,1)$ & 1.00000 & 0.010749346 & 0.010749346 \\
$(2,1)$ & 3.11829 & 0.00013294108 & 0.00013294108 \\
$(1,3)$ & 0.494632 & 0.067442022 & 0.067442022 \\*
\end{longtable}}
The momenta $(1.4302\ldots,2.5)$ have $K=1$ and give the same $|\beta|^2$ as $(1,1)$; an isotropic expansion creates no massless pairs ($|\beta|^2<10^{-20}$). With a mass the two classes $\mathcal Z=\pm i$ are created at different rates: for $m=0.3$, $|\vec p|=1$, $|\vec q|=1.5$, $x_0$ from 0.05 to 40,
\begin{equation*}
|\beta|^2_{16\times16}=0.027052647=\tfrac12\big(|\beta|^2_{+m}+|\beta|^2_{-m}\big),\qquad |\beta|^2_{+m}=0.025286593,\quad |\beta|^2_{-m}=0.028818702 .
\end{equation*}
\Needspace{8\baselineskip}\subsection*{7.3 Modern epoch: a static auxiliary space}\addcontentsline{toc}{subsection}{7.3 Modern epoch: a static auxiliary space}
With $b=b_0$ ($H_b=0$) the Lovelock tensors reduce to (entries 00, 11, 44)
\begin{equation*}
E^{(1)}=\mathrm{diag}\big(3H_a^2,\ 2H_a'+3H_a^2,\ 3H_a'+6H_a^2\big),\qquad E^{(2)}=\mathrm{diag}\big(0,\,0,\,-12H_a^2(H_a'+H_a^2)\big),\qquad E^{(3)}=0 .
\end{equation*}
The 00 and 11 equations are the Friedmann equations of ordinary cosmology,
\begin{equation*}
3H_a^2=\kappa\varepsilon+\Lambda=8\pi G_4\,\rho_4+\Lambda,\qquad 2H_a'+3H_a^2=-\kappa p_a+\Lambda,\qquad G_4=\frac{G_8}{V_\perp},\quad V_\perp=b_0^4V_c,
\end{equation*}
for matter confined to ordinary space ($\varepsilon=\rho_4/V_\perp$; $V_c$ the coordinate volume of the auxiliary space). The 44 equation then FIXES the stress that holds $b$ constant: $\kappa p_b=\Lambda-3H_a'-6H_a^2+12\alpha_2H_a^2(H_a'+H_a^2)$ — the analogue of the fold fluid of Dirac\_claude. A mode with auxiliary momentum $\vec q$ moves in ordinary space as a particle of mass $\sqrt{m^2+|\vec q|^2/b_0^2}$ (a Kaluza–Klein tower on a torus). Pair creation today is suppressed by $e^{-2\pi M/H_0}$ (section 7.4) with $H_0\approx1.4\times10^{-42}$ GeV, i.e. absent.

\Needspace{8\baselineskip}\subsection*{7.4 Far future: de Sitter expansion and thermal pair creation}\addcontentsline{toc}{subsection}{7.4 Far future: de Sitter expansion and thermal pair creation}
Eight-dimensional de Sitter space $a=b=e^{Hx_0}$ is an exact vacuum solution when all $E^{(k)}$ are proportional to $\delta$:
\begin{equation*}
E^{(1)}=21H^2\,\delta,\quad E^{(2)}=-420H^4\,\delta,\quad E^{(3)}=2520H^6\,\delta\qquad\Longrightarrow\qquad 21\alpha_1H^2-420\alpha_2H^4+2520\alpha_3H^6=\Lambda .
\end{equation*}
Over a static auxiliary space ($a=e^{Hx_0}$, $b=b_0$): $E^{(1)}=\mathrm{diag}(3,3,6)H^2$, $E^{(2)}=\mathrm{diag}(0,0,-12)H^4$, $E^{(3)}=0$. In both cases a mode of dirac16complex obeys $i\,dc/dx_0=(k\,e^{-Hx_0}\sigma_1+M\sigma_3)c$, with $M=m$, or with the Kaluza–Klein mass $M=\sqrt{m^2+|\vec q|^2/b_0^2}$. Starting from the adiabatic vacuum at early times the number of created pairs per mode is exactly the Fermi–Dirac occupation at the Gibbons–Hawking temperature $T=H/(2\pi)$, independent of $k$:
\begin{equation*}
|\beta|^2=\frac{1}{e^{2\pi M/H}+1}.
\end{equation*}
{\small
\begin{longtable}{>{\raggedright\arraybackslash}p{\dimexpr 0.2000\linewidth-2\tabcolsep\relax}>{\raggedright\arraybackslash}p{\dimexpr 0.4000\linewidth-2\tabcolsep\relax}>{\raggedright\arraybackslash}p{\dimexpr 0.4000\linewidth-2\tabcolsep\relax}}
\toprule
\textbf{$M/H$} & \textbf{$|\beta|^2$ (Mathematica)} & \textbf{$1/(e^{2\pi M/H}+1)$} \\
\midrule\endfirsthead
\toprule
\textbf{$M/H$} & \textbf{$|\beta|^2$ (Mathematica)} & \textbf{$1/(e^{2\pi M/H}+1)$} \\
\midrule\endhead
\endfoot
\bottomrule\endlastfoot
0.25 & 0.172102919 & 0.172102899 \\
0.5 & 0.0414238845 & 0.0414238322 \\
1.0 & 0.00186397746 & 0.00186396189 \\*
\end{longtable}}
The integrator is also checked on the Landau–Zener problem $i\,dc/dt=(\kappa\sigma_1+t\sigma_3)c$, $t$ from $-150$ to $150$, whose exact result is $|\beta|^2=e^{-\pi\kappa^2}$.

\Needspace{12\baselineskip}\section*{8. Verification}\addcontentsline{toc}{section}{8. Verification}
Every statement above is a check of 7+1.nb (Wolfram Language 15.0.1) and of the independent Python program tools/ref\_gravity.py (sympy, numpy, scipy), which rebuilds the Dirac matrices from 4+4.nb’s split octonions and computes everything again with its own curvature and Lovelock engine (tools/lovelock71.py). The checks of sections 11–14 of the notebook:

{\footnotesize
\begin{longtable}{>{\raggedright\arraybackslash}p{\dimexpr 0.0500\linewidth-2\tabcolsep\relax}>{\raggedright\arraybackslash}p{\dimexpr 0.7500\linewidth-2\tabcolsep\relax}>{\raggedright\arraybackslash}p{\dimexpr 0.0800\linewidth-2\tabcolsep\relax}>{\raggedright\arraybackslash}p{\dimexpr 0.1200\linewidth-2\tabcolsep\relax}}
\toprule
\textbf{\#} & \textbf{check} & \textbf{result} & \textbf{section} \\
\midrule\endfirsthead
\toprule
\textbf{\#} & \textbf{check} & \textbf{result} & \textbf{section} \\
\midrule\endhead
\endfoot
\bottomrule\endlastfoot
1 & \texttt{cosmological metric: sectional curvatures K\_0i = -(Ha'+Ha\textasciicircum{}2), K\_0k = -(Hb'+Hb\textasciicircum{}2), K\_ij = -Ha\textasciicircum{}2, K\_kl = -Hb\textasciicircum{}2, K\_ik = -Ha Hb} & True & 11.2 \\
2 & \texttt{cosmological metric: R\textasciicircum{}cd\_ab is non-zero only for \{c,d\} = \{a,b\}} & True & 11.2 \\
3 & \texttt{Einstein tensor: G\textasciicircum{}0\_0 = 3Ha\textasciicircum{}2 + 12 Ha Hb + 6 Hb\textasciicircum{}2, G\textasciicircum{}1\_1 = 2Ha' + 3Ha\textasciicircum{}2 + 4Hb' + 10Hb\textasciicircum{}2 + 8HaHb, G\textasciicircum{}4\_4 = 3Ha' + 6Ha\textasciicircum{}2 + 3Hb' + 6Hb\textasciicircum{}2 + 9HaHb} & True & 11.2 \\
4 & \texttt{scalar curvature R = -2 (3Ha' + 6Ha\textasciicircum{}2 + 4Hb' + 10Hb\textasciicircum{}2 + 12 Ha Hb)} & True & 11.2 \\
5 & \texttt{inhomogeneous metric, exactly at (x0, x4) = (1/2, 1/3): R\_rsmn = -R\_srmn = -R\_rsnm = R\_mnrs and R\textasciicircum{}r\_[smn] = 0 (first Bianchi identity)} & True & 11.3 \\
6 & \texttt{inhomogeneous metric: contracted Bianchi identity nabla\_mu G\textasciicircum{}mu\_nu == 0, exactly at (x0, x4) = (1/2, 1/3)} & True & 11.3 \\
7 & \texttt{cosmological metric: contracted Bianchi identity nabla\_mu G\textasciicircum{}mu\_nu = 0 (arbitrary a, b)} & True & 11.3 \\
8 & \texttt{generalized Kronecker delta: contraction identity delta\textasciicircum{}\{m1..mp r\}\_\{n1..np r\} = (8-p) delta\textasciicircum{}\{m1..mp\}\_\{n1..np\} (p = 1..4, 60 random index sets each)} & True & 12.1 \\
9 & \texttt{generalized Kronecker delta with 9 upper indices vanishes in 8 dimensions (100 random index sets)} & True & 12.1 \\
10 & \texttt{xAct: Lovelock1 == R\_a\textasciicircum{}b - (1/2) delta\_a\textasciicircum{}b R, the Einstein tensor (abstract indices)} & True & 12.3 \\
11 & \texttt{xAct: trace of Lovelock2 == -2 L\_2 with L\_2 = R\textasciicircum{}2 - 4 R\_cd R\textasciicircum{}cd + R\_cdfi R\textasciicircum{}cdfi (abstract indices)} & True & 12.3 \\
12 & \texttt{xAct, cosmological metric: Lovelock1 == EinsteinCD == the Einstein tensor of the curvature engine of section 11.1} & True & 12.4 \\
13 & \texttt{pair-set algorithm == brute-force generalized-delta contraction for E\textasciicircum{}(1) (cosmological metric)} & True & 12.5 \\
14 & \texttt{pair-set algorithm == xAct (section 12.4) for E\textasciicircum{}(1), E\textasciicircum{}(2), E\textasciicircum{}(3) of the cosmological metric (two independent computations)} & True & 12.5 \\
15 & \texttt{E\textasciicircum{}(1) == Einstein tensor G and L\_1 == R (cosmological metric)} & True & 12.6 \\
16 & \texttt{trace identity E\textasciicircum{}(k)mu\_mu == -((8-2k)/2) L\_k for k = 0..4} & True & 12.6 \\
17 & \texttt{E\textasciicircum{}(4) == 0 identically in 8 dimensions, while the Euler density L\_4 != 0} & True & 12.6 \\
18 & \texttt{every E\textasciicircum{}(k) is diagonal, E\textasciicircum{}(k)1\_1 = E\textasciicircum{}(k)2\_2 = E\textasciicircum{}(k)3\_3 and E\textasciicircum{}(k)4\_4 = ... = E\textasciicircum{}(k)7\_7} & True & 12.6 \\
19 & \texttt{E\textasciicircum{}(2) == Lanczos tensor and L\_2 == R\textasciicircum{}2 - 4 Ric.Ric + Riem.Riem (cosmological metric)} & True & 12.6 \\
20 & \texttt{E\textasciicircum{}(2)0\_0 = -12 Hb (4Ha\textasciicircum{}3 + 18Ha\textasciicircum{}2 Hb + 12 Ha Hb\textasciicircum{}2 + Hb\textasciicircum{}3) and E\textasciicircum{}(3)0\_0 = 360 Ha\textasciicircum{}2 Hb\textasciicircum{}3 (4Ha + 3Hb)} & True & 12.6 \\
21 & \texttt{divergence-free: d0 E\textasciicircum{}0\_0 + 3Ha (E\textasciicircum{}0\_0 - E\textasciicircum{}1\_1) + 4Hb (E\textasciicircum{}0\_0 - E\textasciicircum{}4\_4) == 0 for k = 1, 2, 3 (arbitrary a, b)} & True & 12.6 \\
22 & \texttt{variation of nL a\textasciicircum{}3 b\textasciicircum{}4 L\_k: EL\_nL = -2 a\textasciicircum{}3 b\textasciicircum{}4 E\textasciicircum{}(k)0\_0, EL\_a = -6 a\textasciicircum{}2 b\textasciicircum{}4 E\textasciicircum{}(k)1\_1, EL\_b = -8 a\textasciicircum{}3 b\textasciicircum{}3 E\textasciicircum{}(k)4\_4 (k = 1..4)} & True & 12.6 \\
23 & \texttt{inhomogeneous metric: nabla\_mu E\textasciicircum{}(k)mu\_nu == 0 EXACTLY at (x0, x4) = (1/2, 1/3) for k = 1, 2, 3} & True & 12.7 \\
24 & \texttt{inhomogeneous metric: E\textasciicircum{}(1) == G at the point, and E\textasciicircum{}(k)\_mu nu symmetric (k = 1, 2, 3)} & True & 12.7 \\
25 & \texttt{inhomogeneous metric: E\textasciicircum{}(2) == Lanczos tensor at the point (exact)} & True & 12.7 \\
26 & \texttt{pp-wave: E\textasciicircum{}(1)\_mu nu == +(1/2)(d2\textasciicircum{}2 + d3\textasciicircum{}2) fPP l\_mu l\_nu} & True & 12.8 \\
27 & \texttt{pp-wave: E\textasciicircum{}(2) == E\textasciicircum{}(3) == 0 for every profile fPP} & True & 12.8 \\
28 & \texttt{pp-wave: xAct (section 12.4) == pair-set algorithm for E\textasciicircum{}(1), E\textasciicircum{}(2), E\textasciicircum{}(3) (a metric that is not diagonal; arbitrary profile fPP)} & True & 12.8 \\
29 & \texttt{exponential expansion a = e\textasciicircum{}\{H x0\}, b = e\textasciicircum{}\{qH H x0\}: xAct Lovelock\_k == the closed forms of section 12.6 with Ha = H, Hb = qH H, Ha' = Hb' = 0 (k = 1, 2, 3)} & True & 12.9 \\
30 & \texttt{exponential expansion: three distinct vacuum equations, related by 3 Ha (eq00 - eq11) + 4 Hb (eq00 - eq44) == 0} & True & 12.9 \\
31 & \texttt{exponential vacuum solutions: (w2, w3) solve all 64 components for every qH (generic qH) and every Lambda} & True & 12.9 \\
32 & \texttt{eternal inflation (a = e\textasciicircum{}\{H x0\}) with deflation (b = e\textasciicircum{}\{-H x0\}): w2 == Lambda/4 - w1/2, w3 == w1/24 - Lambda/36; for Lambda = 0: w2/w1 == -1/2, w3/w1 == 1/24} & True & 12.9 \\
33 & \texttt{cosmological metric: Omega\_0 == 0, Omega\_i == (a'/2) Gamma17[i].Gamma17[0], Omega\_k == (b'/2) Gamma17[k].Gamma17[0]} & True & 13.2 \\
34 & \texttt{cosmological metric: Gamma\textasciicircum{}mu Omega\_mu == (3a'/(2a) + 2b'/b) Gamma17[0]} & True & 13.2 \\
35 & \texttt{current identity: d\_mu(sqrt|g| psibar Gamma\textasciicircum{}mu psi) == sqrt|g| [(D\_mu psibar) Gamma\textasciicircum{}mu psi + psibar Gamma\textasciicircum{}mu D\_mu psi] (psi, psibar arbitrary)} & True & 13.2 \\
36 & \texttt{Lichnerowicz formula: (Gamma\textasciicircum{}mu D\_mu)\textasciicircum{}2 psi == (g\textasciicircum{}mu nu nabla\_mu D\_nu - R/4) psi} & True & 13.2 \\
37 & \texttt{rescaling psi = (a\textasciicircum{}3 b\textasciicircum{}4)\textasciicircum{}(-1/2) chi removes the spin connection: Gamma\textasciicircum{}mu D\_mu psi == (a\textasciicircum{}3 b\textasciicircum{}4)\textasciicircum{}(-1/2) Gamma\textasciicircum{}mu d\_mu chi} & True & 13.2 \\
38 & \texttt{mode equation: Gamma17[0].(i Gamma\textasciicircum{}mu D\_mu - m) psi == (a\textasciicircum{}3 b\textasciicircum{}4)\textasciicircum{}(-1/2) e\textasciicircum{}\{i(p.x+q.y)\} (i dc/dx0 - h(x0) c) for psi = (a\textasciicircum{}3 b\textasciicircum{}4)\textasciicircum{}(-1/2) e\textasciicircum{}\{i(p.x+q.y)\} c(x0)} & True & 13.2 \\
39 & \texttt{conformal weight 7/2: Gamma\textasciicircum{}mu D\_mu (C\textasciicircum{}(-7/2) f) == C\textasciicircum{}(-9/2) Gamma17up[A] d\_A f for g = C(x0)\textasciicircum{}2 eta17} & True & 13.4 \\
40 & \texttt{Cartan's first structure equation de\textasciicircum{}A + omega\textasciicircum{}A\_B \textasciicircum{} e\textasciicircum{}B == 0 (no torsion): cosmological, pp-wave and inhomogeneous frames} & True & 13.5 \\
41 & \texttt{structure coefficients: omega\_C,AB == -(1/2)(C\_ACB + C\_BAC - C\_CBA) with C\textasciicircum{}A\_BC = e\_B\textasciicircum{}mu e\_C\textasciicircum{}nu (d\_mu e\textasciicircum{}A\_nu - d\_nu e\textasciicircum{}A\_mu): cosmological and inhomogeneous frames} & True & 13.5 \\
42 & \texttt{static radial metric: Cartan's first structure equation and omega\_C,AB == -(1/2)(C\_ACB + C\_BAC - C\_CBA)} & True & 13.5 \\
43 & \texttt{static radial metric: Omega\_0 == (f'/(4 Sqrt[f] Sqrt[h])) Gamma17[1].Gamma17[0], Omega\_1 == 0, Omega\_k == (1/(2 Sqrt[h])) Gamma17[1].Gamma17[k] (k = 2..7)} & True & 13.5 \\
44 & \texttt{static radial metric: Gamma\textasciicircum{}mu Omega\_mu == (1/Sqrt[h]) (f'/(4 f) + 3/r) Gamma17up[1]} & True & 13.5 \\
45 & \texttt{static radial metric: psi = f\textasciicircum{}(-1/4) r\textasciicircum{}(-3) chi removes the spin connection, Gamma\textasciicircum{}mu D\_mu psi == f\textasciicircum{}(-1/4) r\textasciicircum{}(-3) Gamma\textasciicircum{}mu d\_mu chi} & True & 13.5 \\
46 & \texttt{U(1) gauge covariance: Gamma\textasciicircum{}mu (D\_mu - i q A'\_mu)(e\textasciicircum{}\{i q alpha\} psi) == e\textasciicircum{}\{i q alpha\} Gamma\textasciicircum{}mu (D\_mu - i q A\_mu) psi with A'\_mu = A\_mu + d\_mu alpha} & True & 13.5 \\
47 & \texttt{current identity with the gauge field: d\_mu(sqrt|g| psibar Gamma\textasciicircum{}mu psi) == sqrt|g| [(D\_mu psibar) Gamma\textasciicircum{}mu psi + psibar Gamma\textasciicircum{}mu D\_mu psi], D including -i q A\_mu} & True & 13.5 \\
48 & \texttt{Kasner metric: the Einstein tensor computed directly from the metric equals the closed form (x0\textasciicircum{}2 G\textasciicircum{}mu\_mu, mu = 0, 1, 4)} & True & 14.1 \\
49 & \texttt{Einstein Kasner: (p, q) = ((1+2Sqrt[2])/7, (2-3Sqrt[2])/14) and ((1-2Sqrt[2])/7, (2+3Sqrt[2])/14) solve E\textasciicircum{}(1) == 0} & True & 14.1 \\
50 & \texttt{Gauss-Bonnet Kasner: 3 real roots of the cubic; (p, (3-3p)/4) solve E\textasciicircum{}(2) == 0} & True & 14.1 \\
51 & \texttt{cubic Lovelock Kasner: (p, q) = (-15/7, 20/7) solves E\textasciicircum{}(3) == 0; E\textasciicircum{}(3) == 0 identically for q = 0 and for p = 0} & True & 14.1 \\
52 & \texttt{2x2 reduction: A, B, betaD square to ID16 and anticommute; Zc = A.B.betaD: Zc.Zc == -ID16, Transpose[Zc] == -Zc, Zc commutes with A, B, betaD} & True & 14.2 \\
53 & \texttt{Kasner pair creation: |alpha|\textasciicircum{}2 + |beta|\textasciicircum{}2 == 1 (to 10\textasciicircum{}-8) for (|p|, |q|) = (1,1), (2,1), (1,3)} & True & 14.3 \\
54 & \texttt{Kasner scaling: (|p|, |q|) = (1.4302.., 2.5) has K = 1 and the same |beta|\textasciicircum{}2 as (1, 1) (to 10\textasciicircum{}-6)} & True & 14.3 \\
55 & \texttt{16x16 evolution == 4 copies of the 2x2 evolution with +m sigma3 and 4 with -m sigma3 (Kasner, m = 0.3, |p| = 1, |q| = 1.5, x0 from 0.05 to 40): average |beta|\textasciicircum{}2 agrees to 10\textasciicircum{}-8; unitarity; the two classes differ} & True & 14.3 \\
56 & \texttt{isotropic expansion, massless: no pair creation (|beta|\textasciicircum{}2 < 10\textasciicircum{}-20)} & True & 14.3 \\
57 & \texttt{static auxiliary space (Hb = 0): E\textasciicircum{}(1) = (3Ha\textasciicircum{}2, 2Ha' + 3Ha\textasciicircum{}2, 3Ha' + 6Ha\textasciicircum{}2), E\textasciicircum{}(2) = (0, 0, -12 Ha\textasciicircum{}2 (Ha' + Ha\textasciicircum{}2)), E\textasciicircum{}(3) = 0} & True & 14.4 \\
58 & \texttt{8D de Sitter: E\textasciicircum{}(1) = 21 H\textasciicircum{}2, E\textasciicircum{}(2) = -420 H\textasciicircum{}4, E\textasciicircum{}(3) = 2520 H\textasciicircum{}6 (all diagonal entries)} & True & 14.5 \\
59 & \texttt{de Sitter over a static auxiliary space: E\textasciicircum{}(1) = (3, 3, 6) H\textasciicircum{}2, E\textasciicircum{}(2) = (0, 0, -12) H\textasciicircum{}4, E\textasciicircum{}(3) = 0} & True & 14.5 \\
60 & \texttt{de Sitter pair creation: |beta|\textasciicircum{}2 == 1/(Exp[2 Pi M/H] + 1) to 10\textasciicircum{}-7, M/H = 0.25, 0.5, 1} & True & 14.5 \\
61 & \texttt{integrator benchmark (Landau-Zener): |beta|\textasciicircum{}2 == Exp[-Pi kappa\textasciicircum{}2] to 10\textasciicircum{}-6} & True & 14.5 \\
62 & \texttt{cosmological metric: eta17 e e == g and \{Gamma\_mu, Gamma\_nu\} == 2 g\_mu nu ID16 (generalized frame matrices, real)} & True & 13.1 \\
63 & \texttt{cosmological metric: canonical spin connection antisymmetric, omega\_mu AB == -omega\_mu BA} & True & 13.1 \\
64 & \texttt{cosmological metric: vielbein postulate d\_mu e\textasciicircum{}A\_nu - \{l mu nu\} e\textasciicircum{}A\_l + omega\_mu\textasciicircum{}A\_B e\textasciicircum{}B\_nu == 0} & True & 13.1 \\
65 & \texttt{cosmological metric: TOTAL covariant derivative of the frame matrices vanishes, d\_mu Gamma\_nu - \{l mu nu\} Gamma\_l + [Omega\_mu, Gamma\_nu] == 0} & True & 13.1 \\
66 & \texttt{cosmological metric: constant frame matrices are covariantly constant, omega\_mu\textasciicircum{}A\_B Gamma17up[B] + [Omega\_mu, Gamma17up[A]] == 0} & True & 13.1 \\
67 & \texttt{cosmological metric: Omega\_mu real and Gamma17[0].Transpose[Omega\_mu].Gamma17[0] == -Omega\_mu (so D\_mu psibar = d\_mu psibar - psibar Omega\_mu)} & True & 13.1 \\
68 & \texttt{cosmological metric: spin curvature F\_mu nu == (1/4) R\_AB mu nu Gamma17up[A].Gamma17up[B]} & True & 13.1 \\
69 & \texttt{cosmological metric: Bianchi identity d\_l F\_mn + [Omega\_l, F\_mn] + cyclic == 0} & True & 13.1 \\
70 & \texttt{cosmological metric: (1/2) Gamma\textasciicircum{}mu Gamma\textasciicircum{}nu F\_mu nu == -(R/4) ID16} & True & 13.1 \\
71 & \texttt{pp-wave: eta17 e e == g and \{Gamma\_mu, Gamma\_nu\} == 2 g\_mu nu ID16 (generalized frame matrices, real)} & True & 13.1 \\
72 & \texttt{pp-wave: canonical spin connection antisymmetric, omega\_mu AB == -omega\_mu BA} & True & 13.1 \\
73 & \texttt{pp-wave: vielbein postulate d\_mu e\textasciicircum{}A\_nu - \{l mu nu\} e\textasciicircum{}A\_l + omega\_mu\textasciicircum{}A\_B e\textasciicircum{}B\_nu == 0} & True & 13.1 \\
74 & \texttt{pp-wave: TOTAL covariant derivative of the frame matrices vanishes, d\_mu Gamma\_nu - \{l mu nu\} Gamma\_l + [Omega\_mu, Gamma\_nu] == 0} & True & 13.1 \\
75 & \texttt{pp-wave: constant frame matrices are covariantly constant, omega\_mu\textasciicircum{}A\_B Gamma17up[B] + [Omega\_mu, Gamma17up[A]] == 0} & True & 13.1 \\
76 & \texttt{pp-wave: Omega\_mu real and Gamma17[0].Transpose[Omega\_mu].Gamma17[0] == -Omega\_mu (so D\_mu psibar = d\_mu psibar - psibar Omega\_mu)} & True & 13.1 \\
77 & \texttt{pp-wave: spin curvature F\_mu nu == (1/4) R\_AB mu nu Gamma17up[A].Gamma17up[B]} & True & 13.1 \\
78 & \texttt{pp-wave: Bianchi identity d\_l F\_mn + [Omega\_l, F\_mn] + cyclic == 0} & True & 13.1 \\
79 & \texttt{pp-wave: (1/2) Gamma\textasciicircum{}mu Gamma\textasciicircum{}nu F\_mu nu == -(R/4) ID16} & True & 13.1 \\
80 & \texttt{inhomogeneous metric (exact at (x0, x4) = (1/2, 1/3)): eta17 e e == g and \{Gamma\_mu, Gamma\_nu\} == 2 g\_mu nu ID16 (generalized frame matrices, real)} & True & 13.1 \\
81 & \texttt{inhomogeneous metric (exact at (x0, x4) = (1/2, 1/3)): canonical spin connection antisymmetric, omega\_mu AB == -omega\_mu BA} & True & 13.1 \\
82 & \texttt{inhomogeneous metric (exact at (x0, x4) = (1/2, 1/3)): vielbein postulate d\_mu e\textasciicircum{}A\_nu - \{l mu nu\} e\textasciicircum{}A\_l + omega\_mu\textasciicircum{}A\_B e\textasciicircum{}B\_nu == 0} & True & 13.1 \\
83 & \texttt{inhomogeneous metric (exact at (x0, x4) = (1/2, 1/3)): TOTAL covariant derivative of the frame matrices vanishes, d\_mu Gamma\_nu - \{l mu nu\} Gamma\_l + [Omega\_mu, Gamma\_nu] == 0} & True & 13.1 \\
84 & \texttt{inhomogeneous metric (exact at (x0, x4) = (1/2, 1/3)): constant frame matrices are covariantly constant, omega\_mu\textasciicircum{}A\_B Gamma17up[B] + [Omega\_mu, Gamma17up[A]] == 0} & True & 13.1 \\
85 & \texttt{inhomogeneous metric (exact at (x0, x4) = (1/2, 1/3)): Omega\_mu real and Gamma17[0].Transpose[Omega\_mu].Gamma17[0] == -Omega\_mu (so D\_mu psibar = d\_mu psibar - psibar Omega\_mu)} & True & 13.1 \\
86 & \texttt{inhomogeneous metric (exact at (x0, x4) = (1/2, 1/3)): spin curvature F\_mu nu == (1/4) R\_AB mu nu Gamma17up[A].Gamma17up[B]} & True & 13.1 \\
87 & \texttt{inhomogeneous metric (exact at (x0, x4) = (1/2, 1/3)): Bianchi identity d\_l F\_mn + [Omega\_l, F\_mn] + cyclic == 0} & True & 13.1 \\
88 & \texttt{inhomogeneous metric (exact at (x0, x4) = (1/2, 1/3)): (1/2) Gamma\textasciicircum{}mu Gamma\textasciicircum{}nu F\_mu nu == -(R/4) ID16} & True & 13.1 \\
89 & \texttt{static radial metric: eta17 e e == g and \{Gamma\_mu, Gamma\_nu\} == 2 g\_mu nu ID16 (generalized frame matrices, real)} & True & 13.1 \\
90 & \texttt{static radial metric: canonical spin connection antisymmetric, omega\_mu AB == -omega\_mu BA} & True & 13.1 \\
91 & \texttt{static radial metric: vielbein postulate d\_mu e\textasciicircum{}A\_nu - \{l mu nu\} e\textasciicircum{}A\_l + omega\_mu\textasciicircum{}A\_B e\textasciicircum{}B\_nu == 0} & True & 13.1 \\
92 & \texttt{static radial metric: TOTAL covariant derivative of the frame matrices vanishes, d\_mu Gamma\_nu - \{l mu nu\} Gamma\_l + [Omega\_mu, Gamma\_nu] == 0} & True & 13.1 \\
93 & \texttt{static radial metric: constant frame matrices are covariantly constant, omega\_mu\textasciicircum{}A\_B Gamma17up[B] + [Omega\_mu, Gamma17up[A]] == 0} & True & 13.1 \\
94 & \texttt{static radial metric: Omega\_mu real and Gamma17[0].Transpose[Omega\_mu].Gamma17[0] == -Omega\_mu (so D\_mu psibar = d\_mu psibar - psibar Omega\_mu)} & True & 13.1 \\
95 & \texttt{static radial metric: spin curvature F\_mu nu == (1/4) R\_AB mu nu Gamma17up[A].Gamma17up[B]} & True & 13.1 \\
96 & \texttt{static radial metric: Bianchi identity d\_l F\_mn + [Omega\_l, F\_mn] + cyclic == 0} & True & 13.1 \\
97 & \texttt{static radial metric: (1/2) Gamma\textasciicircum{}mu Gamma\textasciicircum{}nu F\_mu nu == -(R/4) ID16} & True & 13.1 \\*
\end{longtable}}
The Python program’s boolean checks:

{\footnotesize
\begin{longtable}{>{\raggedright\arraybackslash}p{\dimexpr 0.0500\linewidth-2\tabcolsep\relax}>{\raggedright\arraybackslash}p{\dimexpr 0.8300\linewidth-2\tabcolsep\relax}>{\raggedright\arraybackslash}p{\dimexpr 0.1200\linewidth-2\tabcolsep\relax}}
\toprule
\textbf{\#} & \textbf{check} & \textbf{result} \\
\midrule\endfirsthead
\toprule
\textbf{\#} & \textbf{check} & \textbf{result} \\
\midrule\endhead
\endfoot
\bottomrule\endlastfoot
1 & \texttt{cosmological metric: sectional curvatures R\textasciicircum{}\{AB\}\_\{AB\} = -(Ha'+Ha\textasciicircum{}2), -(Hb'+Hb\textasciicircum{}2), -Ha\textasciicircum{}2, -Hb\textasciicircum{}2, -Ha Hb} & True \\
2 & \texttt{cosmological metric: Riemann R\textasciicircum{}\{cd\}\_\{ab\} non-zero only for \{c,d\} = \{a,b\}} & True \\
3 & \texttt{G\textasciicircum{}0\_0 = 3Ha\textasciicircum{}2 + 12 Ha Hb + 6 Hb\textasciicircum{}2} & True \\
4 & \texttt{G\textasciicircum{}1\_1 = 2Ha' + 3Ha\textasciicircum{}2 + 4Hb' + 10Hb\textasciicircum{}2 + 8HaHb} & True \\
5 & \texttt{G\textasciicircum{}4\_4 = 3Ha' + 6Ha\textasciicircum{}2 + 3Hb' + 6Hb\textasciicircum{}2 + 9HaHb} & True \\
6 & \texttt{scalar curvature R = -2(3Ha' + 6Ha\textasciicircum{}2 + 4Hb' + 10Hb\textasciicircum{}2 + 12 HaHb)} & True \\
7 & \texttt{E\textasciicircum{}(1) = Einstein tensor G (cosmological metric)} & True \\
8 & \texttt{L\_1 = R} & True \\
9 & \texttt{E\textasciicircum{}(0) = -delta/2} & True \\
10 & \texttt{trace identity E\textasciicircum{}(k)mu\_mu = -(8-2k)/2 L\_k, k = 0..4} & True \\
11 & \texttt{E\textasciicircum{}(4) = 0 identically in 8 dimensions, but L\_4 (Euler density) is not 0} & True \\
12 & \texttt{E\textasciicircum{}(k) diagonal: E\textasciicircum{}(k)mu\_nu = 0 for mu != nu, and E\textasciicircum{}(k)1\_1 = E\textasciicircum{}(k)2\_2 = E\textasciicircum{}(k)3\_3, E\textasciicircum{}(k)4\_4 = ... = E\textasciicircum{}(k)7\_7} & True \\
13 & \texttt{E\textasciicircum{}(2)0\_0 = -12 Hb (4Ha\textasciicircum{}3 + 18Ha\textasciicircum{}2Hb + 12HaHb\textasciicircum{}2 + Hb\textasciicircum{}3)} & True \\
14 & \texttt{E\textasciicircum{}(3)0\_0 = 360 Ha\textasciicircum{}2 Hb\textasciicircum{}3 (4Ha + 3Hb)} & True \\
15 & \texttt{E\textasciicircum{}(2) = Lanczos tensor 2(R R\_mn - 2R\_ma R\textasciicircum{}a\_n - 2R\textasciicircum{}ab R\_manb + R\_m\textasciicircum{}abc R\_nabc) - g\_mn L\_2/2, and L\_2 = R\textasciicircum{}2 - 4 Ric.Ric + Riem.Riem (cosmological metric)} & True \\
16 & \texttt{divergence-free E\textasciicircum{}(k), k = 1, 2, 3 (cosmological metric, arbitrary a, b)} & True \\
17 & \texttt{variation of sqrt|g| L\_k (lapse N, a, b): EL\_N = -2 a\textasciicircum{}3 b\textasciicircum{}4 E\textasciicircum{}(k)0\_0, EL\_a = -6 a\textasciicircum{}2 b\textasciicircum{}4 E\textasciicircum{}(k)1\_1, EL\_b = -8 a\textasciicircum{}3 b\textasciicircum{}3 E\textasciicircum{}(k)4\_4, k = 1..4} & True \\
18 & \texttt{inhomogeneous metric (N, A, B of x0 and x4): E\textasciicircum{}(1) = G at the point} & True \\
19 & \texttt{inhomogeneous metric: nabla\_mu E\textasciicircum{}(k)mu\_nu = 0 exactly at x0 = 1/2, x4 = 1/3, k = 1, 2, 3} & True \\
20 & \texttt{inhomogeneous metric: E\textasciicircum{}(k)\_mu nu symmetric, k = 1, 2, 3} & True \\
21 & \texttt{pp-wave: E\textasciicircum{}(1)\_mu nu = +(1/2) (d2\textasciicircum{}2 + d3\textasciicircum{}2) F l\_mu l\_nu, l = dx0 - dx1} & True \\
22 & \texttt{pp-wave: E\textasciicircum{}(2) = E\textasciicircum{}(3) = 0 for every profile F} & True \\
23 & \texttt{Einstein Kasner: p = (1+-2 sqrt2)/7, q = (2-+3 sqrt2)/14 solve E\textasciicircum{}(1) = 0; 3p+4q = 1, 3p\textasciicircum{}2+4q\textasciicircum{}2 = 1} & True \\
24 & \texttt{Gauss-Bonnet Kasner: the three (p, (3-3p)/4) solve E\textasciicircum{}(2) = 0, sum of the 7 exponents = 3} & True \\
25 & \texttt{cubic Lovelock Kasner: (p, q) = (-15/7, 20/7) solves E\textasciicircum{}(3) = 0, 3p + 4q = 5; degenerate families q = 0 and p = 0} & True \\
26 & \texttt{Einstein Kasner is not a cubic-Lovelock vacuum (E\textasciicircum{}(3) != 0)} & True \\
27 & \texttt{8D de Sitter (a = b = e\textasciicircum{}\{H x0\}): E\textasciicircum{}(1) = 21 H\textasciicircum{}2, E\textasciicircum{}(2) = -420 H\textasciicircum{}4, E\textasciicircum{}(3) = 2520 H\textasciicircum{}6 (times delta)} & True \\
28 & \texttt{4D de Sitter on a static T\textasciicircum{}4 (a = e\textasciicircum{}\{H x0\}, b = b0): E\textasciicircum{}(1) = diag(3,3,6) H\textasciicircum{}2, E\textasciicircum{}(2) = diag(0,0,-12) H\textasciicircum{}4, E\textasciicircum{}(3) = 0} & True \\
29 & \texttt{static T\textasciicircum{}4 (Hb = 0): E\textasciicircum{}(1) = (3Ha\textasciicircum{}2, 2Ha'+3Ha\textasciicircum{}2, 3Ha'+6Ha\textasciicircum{}2), E\textasciicircum{}(2) = (0, 0, -12 Ha\textasciicircum{}2 (Ha'+Ha\textasciicircum{}2)), E\textasciicircum{}(3) = 0} & True \\
30 & \texttt{exponential expansion a = e\textasciicircum{}\{H x0\}, b = e\textasciicircum{}\{q H x0\}: E\textasciicircum{}(k) (00, 11, 44) = closed forms with Ha = H, Hb = q H, Ha' = Hb' = 0 (k = 1, 2, 3)} & True \\
31 & \texttt{exponential expansion: 3 Ha (eq00 - eq11) + 4 Hb (eq00 - eq44) = 0 identically (the three vacuum equations have rank 2)} & True \\
32 & \texttt{exponential vacuum solutions: one solution (w2, w3) for given w1, Lambda, q, satisfying all three equations} & True \\
33 & \texttt{eternal inflation a = e\textasciicircum{}\{H x0\} with deflation b = e\textasciicircum{}\{-H x0\}: w2 = Lambda/4 - w1/2, w3 = w1/24 - Lambda/36 (Lambda = 0: w2/w1 = -1/2, w3/w1 = 1/24)} & True \\
34 & \texttt{cosmological metric: eta e e = g} & True \\
35 & \texttt{cosmological metric: \{Gamma\_mu, Gamma\_nu\} = 2 g\_mu nu} & True \\
36 & \texttt{cosmological metric: Gamma\_mu real} & True \\
37 & \texttt{cosmological metric: omega\_mu AB = -omega\_mu BA} & True \\
38 & \texttt{cosmological metric: vielbein postulate d\_mu e\textasciicircum{}A\_nu - chr\textasciicircum{}l\_mu nu e\textasciicircum{}A\_l + omega\_mu\textasciicircum{}A\_B e\textasciicircum{}B\_nu = 0} & True \\
39 & \texttt{cosmological metric: Omega\_mu real, Gamma17[0] Omega\_mu\textasciicircum{}T Gamma17[0] = -Omega\_mu} & True \\
40 & \texttt{cosmological metric: total covariant derivative of the frame matrices: d\_mu Gamma\_nu - chr\textasciicircum{}l\_mu nu Gamma\_l + [Omega\_mu, Gamma\_nu] = 0} & True \\
41 & \texttt{cosmological metric: constant frame matrices: omega\_mu\textasciicircum{}A\_B Gamma\textasciicircum{}B + [Omega\_mu, Gamma\textasciicircum{}A] = 0} & True \\
42 & \texttt{cosmological metric: spin curvature F\_mu nu = (1/4) R\_AB mu nu Gamma\textasciicircum{}A Gamma\textasciicircum{}B} & True \\
43 & \texttt{cosmological metric: Bianchi identity d\_l F\_mn + [Omega\_l, F\_mn] + cyclic = 0} & True \\
44 & \texttt{cosmological metric: (1/2) Gamma\textasciicircum{}mu Gamma\textasciicircum{}nu F\_mu nu = -(R/4) 1\_16} & True \\
45 & \texttt{cosmological metric: Gamma\textasciicircum{}mu Omega\_mu = (3a'/(2a) + 2b'/b) Gamma17[0]} & True \\
46 & \texttt{cosmological metric: Omega\_i = (a'/2) Gamma17[i] Gamma17[0] (i = 1..3), Omega\_k = (b'/2) Gamma17[k] Gamma17[0] (k = 4..7), Omega\_0 = 0} & True \\
47 & \texttt{pp-wave: eta e e = g} & True \\
48 & \texttt{pp-wave: \{Gamma\_mu, Gamma\_nu\} = 2 g\_mu nu} & True \\
49 & \texttt{pp-wave: Gamma\_mu real} & True \\
50 & \texttt{pp-wave: omega\_mu AB = -omega\_mu BA} & True \\
51 & \texttt{pp-wave: vielbein postulate d\_mu e\textasciicircum{}A\_nu - chr\textasciicircum{}l\_mu nu e\textasciicircum{}A\_l + omega\_mu\textasciicircum{}A\_B e\textasciicircum{}B\_nu = 0} & True \\
52 & \texttt{pp-wave: Omega\_mu real, Gamma17[0] Omega\_mu\textasciicircum{}T Gamma17[0] = -Omega\_mu} & True \\
53 & \texttt{pp-wave: total covariant derivative of the frame matrices: d\_mu Gamma\_nu - chr\textasciicircum{}l\_mu nu Gamma\_l + [Omega\_mu, Gamma\_nu] = 0} & True \\
54 & \texttt{pp-wave: constant frame matrices: omega\_mu\textasciicircum{}A\_B Gamma\textasciicircum{}B + [Omega\_mu, Gamma\textasciicircum{}A] = 0} & True \\
55 & \texttt{pp-wave: spin curvature F\_mu nu = (1/4) R\_AB mu nu Gamma\textasciicircum{}A Gamma\textasciicircum{}B} & True \\
56 & \texttt{pp-wave: Bianchi identity d\_l F\_mn + [Omega\_l, F\_mn] + cyclic = 0} & True \\
57 & \texttt{pp-wave: (1/2) Gamma\textasciicircum{}mu Gamma\textasciicircum{}nu F\_mu nu = -(R/4) 1\_16} & True \\
58 & \texttt{inhomogeneous metric N, A, B (x0, x4): eta e e = g} & True \\
59 & \texttt{inhomogeneous metric N, A, B (x0, x4): \{Gamma\_mu, Gamma\_nu\} = 2 g\_mu nu} & True \\
60 & \texttt{inhomogeneous metric N, A, B (x0, x4): Gamma\_mu real} & True \\
61 & \texttt{inhomogeneous metric N, A, B (x0, x4): omega\_mu AB = -omega\_mu BA} & True \\
62 & \texttt{inhomogeneous metric N, A, B (x0, x4): vielbein postulate d\_mu e\textasciicircum{}A\_nu - chr\textasciicircum{}l\_mu nu e\textasciicircum{}A\_l + omega\_mu\textasciicircum{}A\_B e\textasciicircum{}B\_nu = 0} & True \\
63 & \texttt{inhomogeneous metric N, A, B (x0, x4): Omega\_mu real, Gamma17[0] Omega\_mu\textasciicircum{}T Gamma17[0] = -Omega\_mu} & True \\
64 & \texttt{inhomogeneous metric N, A, B (x0, x4): total covariant derivative of the frame matrices: d\_mu Gamma\_nu - chr\textasciicircum{}l\_mu nu Gamma\_l + [Omega\_mu, Gamma\_nu] = 0} & True \\
65 & \texttt{inhomogeneous metric N, A, B (x0, x4): constant frame matrices: omega\_mu\textasciicircum{}A\_B Gamma\textasciicircum{}B + [Omega\_mu, Gamma\textasciicircum{}A] = 0} & True \\
66 & \texttt{inhomogeneous metric N, A, B (x0, x4): spin curvature F\_mu nu = (1/4) R\_AB mu nu Gamma\textasciicircum{}A Gamma\textasciicircum{}B} & True \\
67 & \texttt{inhomogeneous metric N, A, B (x0, x4): Bianchi identity d\_l F\_mn + [Omega\_l, F\_mn] + cyclic = 0} & True \\
68 & \texttt{inhomogeneous metric N, A, B (x0, x4): (1/2) Gamma\textasciicircum{}mu Gamma\textasciicircum{}nu F\_mu nu = -(R/4) 1\_16} & True \\
69 & \texttt{Cartan's first structure equation de\textasciicircum{}A + omega\textasciicircum{}A\_B \textasciicircum{} e\textasciicircum{}B = 0 (torsion-free) for the cosmological frame} & True \\
70 & \texttt{static radial metric f(r) dx0\textasciicircum{}2 - h(r) dr\textasciicircum{}2 - r\textasciicircum{}2 (dx2\textasciicircum{}2 + ... + dx7\textasciicircum{}2): eta e e = g} & True \\
71 & \texttt{static radial metric f(r) dx0\textasciicircum{}2 - h(r) dr\textasciicircum{}2 - r\textasciicircum{}2 (dx2\textasciicircum{}2 + ... + dx7\textasciicircum{}2): \{Gamma\_mu, Gamma\_nu\} = 2 g\_mu nu} & True \\
72 & \texttt{static radial metric f(r) dx0\textasciicircum{}2 - h(r) dr\textasciicircum{}2 - r\textasciicircum{}2 (dx2\textasciicircum{}2 + ... + dx7\textasciicircum{}2): Gamma\_mu real} & True \\
73 & \texttt{static radial metric f(r) dx0\textasciicircum{}2 - h(r) dr\textasciicircum{}2 - r\textasciicircum{}2 (dx2\textasciicircum{}2 + ... + dx7\textasciicircum{}2): omega\_mu AB = -omega\_mu BA} & True \\
74 & \texttt{static radial metric f(r) dx0\textasciicircum{}2 - h(r) dr\textasciicircum{}2 - r\textasciicircum{}2 (dx2\textasciicircum{}2 + ... + dx7\textasciicircum{}2): vielbein postulate d\_mu e\textasciicircum{}A\_nu - chr\textasciicircum{}l\_mu nu e\textasciicircum{}A\_l + omega\_mu\textasciicircum{}A\_B e\textasciicircum{}B\_nu = 0} & True \\
75 & \texttt{static radial metric f(r) dx0\textasciicircum{}2 - h(r) dr\textasciicircum{}2 - r\textasciicircum{}2 (dx2\textasciicircum{}2 + ... + dx7\textasciicircum{}2): Omega\_mu real, Gamma17[0] Omega\_mu\textasciicircum{}T Gamma17[0] = -Omega\_mu} & True \\
76 & \texttt{static radial metric f(r) dx0\textasciicircum{}2 - h(r) dr\textasciicircum{}2 - r\textasciicircum{}2 (dx2\textasciicircum{}2 + ... + dx7\textasciicircum{}2): total covariant derivative of the frame matrices: d\_mu Gamma\_nu - chr\textasciicircum{}l\_mu nu Gamma\_l + [Omega\_mu, Gamma\_nu] = 0} & True \\
77 & \texttt{static radial metric f(r) dx0\textasciicircum{}2 - h(r) dr\textasciicircum{}2 - r\textasciicircum{}2 (dx2\textasciicircum{}2 + ... + dx7\textasciicircum{}2): constant frame matrices: omega\_mu\textasciicircum{}A\_B Gamma\textasciicircum{}B + [Omega\_mu, Gamma\textasciicircum{}A] = 0} & True \\
78 & \texttt{static radial metric f(r) dx0\textasciicircum{}2 - h(r) dr\textasciicircum{}2 - r\textasciicircum{}2 (dx2\textasciicircum{}2 + ... + dx7\textasciicircum{}2): spin curvature F\_mu nu = (1/4) R\_AB mu nu Gamma\textasciicircum{}A Gamma\textasciicircum{}B} & True \\
79 & \texttt{static radial metric f(r) dx0\textasciicircum{}2 - h(r) dr\textasciicircum{}2 - r\textasciicircum{}2 (dx2\textasciicircum{}2 + ... + dx7\textasciicircum{}2): Bianchi identity d\_l F\_mn + [Omega\_l, F\_mn] + cyclic = 0} & True \\
80 & \texttt{static radial metric f(r) dx0\textasciicircum{}2 - h(r) dr\textasciicircum{}2 - r\textasciicircum{}2 (dx2\textasciicircum{}2 + ... + dx7\textasciicircum{}2): (1/2) Gamma\textasciicircum{}mu Gamma\textasciicircum{}nu F\_mu nu = -(R/4) 1\_16} & True \\
81 & \texttt{Cartan's first structure equation for the static radial metric} & True \\
82 & \texttt{structure coefficients: de\textasciicircum{}A = (1/2) C\textasciicircum{}A\_BC e\textasciicircum{}B \textasciicircum{} e\textasciicircum{}C and omega\_C,AB (form index C) = -(1/2)(C\_ACB + C\_BAC - C\_CBA) (static radial metric)} & True \\
83 & \texttt{static radial metric: Omega\_0 = (f'/(4 sqrt(f h))) Gamma17[1] Gamma17[0], Omega\_1 = 0, Omega\_k = (1/(2 sqrt h)) Gamma17[1] Gamma17[k] (k = 2..7)} & True \\
84 & \texttt{static radial metric: Gamma\textasciicircum{}mu Omega\_mu = (1/sqrt h)(f'/(4f) + 3/r) Gamma17up[1], so (i Gamma\textasciicircum{}mu D\_mu - m) psi = [i Gamma17up[0] d0/sqrt f + i Gamma17up[1] (d\_r + f'/(4f) + 3/r)/sqrt h + sum\_k i Gamma17up[k] d\_k/r - m] psi} & True \\
85 & \texttt{static radial metric: psi = f\textasciicircum{}(-1/4) r\textasciicircum{}(-3) chi removes the spin connection: Gamma\textasciicircum{}mu D\_mu psi = f\textasciicircum{}(-1/4) r\textasciicircum{}(-3) Gamma\textasciicircum{}mu d\_mu chi} & True \\
86 & \texttt{U(1) gauge covariance: Gamma\textasciicircum{}mu (D\_mu - i q A'\_mu)(e\textasciicircum{}\{i q alpha\} psi) = e\textasciicircum{}\{i q alpha\} Gamma\textasciicircum{}mu (D\_mu - i q A\_mu) psi with A' = A + d alpha} & True \\
87 & \texttt{current identity with the gauge field: d\_mu(sqrt|g| psibar Gamma\textasciicircum{}mu psi) = sqrt|g| [(D\_mu psibar) Gamma\textasciicircum{}mu psi + psibar Gamma\textasciicircum{}mu D\_mu psi], D including -i q A\_mu} & True \\
88 & \texttt{current identity d\_mu(sqrt|g| psibar Gamma\textasciicircum{}mu psi) = sqrt|g| [(D\_mu psibar) Gamma\textasciicircum{}mu psi + psibar Gamma\textasciicircum{}mu D\_mu psi]} & True \\
89 & \texttt{Lichnerowicz: (Gamma\textasciicircum{}mu D\_mu)\textasciicircum{}2 psi = (g\textasciicircum{}mu nu nabla\_mu D\_nu - R/4) psi} & True \\
90 & \texttt{rescaling psi = (a\textasciicircum{}3 b\textasciicircum{}4)\textasciicircum{}(-1/2) chi removes the spin connection: Gamma\textasciicircum{}mu D\_mu psi = (a\textasciicircum{}3 b\textasciicircum{}4)\textasciicircum{}(-1/2) Gamma\textasciicircum{}mu d\_mu chi} & True \\
91 & \texttt{mode equation: (i Gamma\textasciicircum{}mu D\_mu - m) psi = 0  <=>  i dc/dx0 = h(x0) c, h = alpha.(p/a, q/b) + beta m} & True \\
92 & \texttt{conformal weight 7/2: Gamma\textasciicircum{}mu D\_mu (C\textasciicircum{}(-7/2) f) = C\textasciicircum{}(-9/2) Gamma\textasciicircum{}A d\_A f for g = C(x0)\textasciicircum{}2 eta} & True \\
93 & \texttt{2x2 reduction: A = alpha.n\_a, B = alpha.n\_b, beta: squares 1, mutually anticommuting; Z = A B beta: Z\textasciicircum{}2 = -1, Z\textasciicircum{}T = -Z, commutes with A, B, beta} & True \\
94 & \texttt{Kasner: |alpha|\textasciicircum{}2 + |beta|\textasciicircum{}2 = 1 (to 1e-8)} & True \\
95 & \texttt{Kasner scaling: (|p|, |q|) = (1.430..., 2.5) has K = 1 and the same |beta|\textasciicircum{}2 as (1, 1) (to 1e-6)} & True \\
96 & \texttt{16x16 evolution = 4 copies of the 2x2 evolution with +m sigma3 and 4 copies with -m sigma3: average agrees (to 1e-8), unitarity; the two sectors differ} & True \\
97 & \texttt{isotropic expansion, massless: no pair creation (|beta|\textasciicircum{}2 < 1e-20)} & True \\
98 & \texttt{de Sitter: |beta|\textasciicircum{}2 = 1/(exp(2 pi m/H) + 1) (Fermi-Dirac at the Gibbons-Hawking temperature H/2pi), to 1e-7} & True \\
99 & \texttt{integrator benchmark: Landau-Zener |beta|\textasciicircum{}2 = exp(-pi kappa\textasciicircum{}2) (to 1e-6)} & True \\*
\end{longtable}}
\Needspace{12\baselineskip}\section*{9. How to reproduce}\addcontentsline{toc}{section}{9. How to reproduce}
\begin{enumerate}
\item \texttt{python3 tools/make\_7p1\_nb.py} writes 7+1.nb (unevaluated) from tools/make\_7p1\_nb.py and tools/nb\_part\_*.py.
\item Evaluate 7+1.nb in a fresh kernel of Wolfram 15.0.1 (Evaluation \texorpdfstring{\ensuremath{\triangleright}}{▸} Quit Kernel, then Evaluation \texorpdfstring{\ensuremath{\triangleright}}{▸} Evaluate Notebook; or \texttt{Get["D:/Developer/github/7+1/tools/run\_eval.wl"]}) and save it; the last section summarizes all checks.
\item \texttt{python3 tools/ref\_gravity.py} repeats the checks of this document independently (exit code 0 if all pass).
\item \texttt{python3 tools/make\_docs.py} regenerates this document, its symbol dictionary and its provenance record.
\end{enumerate}
\end{document}
